\clearpage
\phantomsection
\section*{Chương 11. Bộ lọc trung vị Median-3 thay cho bộ dò đường Path Finder}
\addcontentsline{toc}{section}{Chương 11. Bộ lọc trung vị Median-3 thay cho bộ dò đường Path Finder}

\begin{learningobj}
\begin{itemize}
    \item Thấu hiểu bản chất thuật toán quy hoạch động Viterbi trong bộ dò đường (\emph{Path Finder}) của Praat 7.0.03: từ ma trận chi phí cục bộ $\delta$ và ma trận con trỏ truy vết $\psi$ kích thước $N \times 15$ đến các hàm phạt nhảy quãng tám ($0{,}35$) và đổi trạng thái vô thanh / hữu thanh ($0{,}14$).
    \item Phân tích nguyên nhân sâu xa vì sao kiến trúc phần cứng FPGA loại bỏ bộ dò đường Viterbi: nguy cơ bùng nổ bộ nhớ SRAM nội chip và độ trễ toàn tệp phi nhân quả trong luồng xử lý dòng thời gian thực.
    \item Chứng minh toán học và thực nghiệm sự suy biến của Viterbi thành phép tìm cực đại độc lập từng khung (\emph{per-frame argmax}) khi chi phí chuyển trạng thái triệt tiêu về 0; thiết lập cấu hình Praat R2 làm chuẩn đối sánh ngang hàng (\emph{apples-to-apples}) cho ngõ ra FPGA trước bộ lọc trung vị.
    \item Đánh giá định lượng hiệu quả và giới hạn của bộ lọc trung vị 3 điểm (Median-3): triệt tiêu hoàn toàn gai đơn lẻ 1 khung (giảm GPE từ 22 xuống 15 khung trên tệp nhiễu $5\text{ dB}$), bất lực trước lỗi kéo dài 2 khung, và tác động làm mượt quyết định vô thanh (VDE) thông qua sự tham gia của giá trị $0$.
    \item Làm chủ vi kiến trúc nhân quả phần cứng: xuất khung $k-1$ tại chu kỳ nhận khung $k$, cơ chế lặp mẫu biên đối xứng và nhịp xả tự động (\emph{Flush beat}) bảo đảm khớp từng bit với mô hình phần mềm.
    \item Khảo sát chi tiết module Verilog \texttt{ch11\_median3.v} và mô hình song sinh Python Twin, kiểm chứng trên 1.065 chu kỳ mô phỏng với 609 kết quả xuất hoàn chỉnh, cùng mã tham chiếu C trên bộ xử lý HPS.
\end{itemize}
\end{learningobj}

% ============================================================================
% 11.1. BỘ DÒ ĐƯỜNG PATH FINDER CỦA PRAAT 7.0.03
% ============================================================================
\subsection*{11.1. Bộ dò đường Path Finder của Praat 7.0.03 và nguyên nhân loại bỏ Viterbi trên FPGA}
\addcontentsline{toc}{subsection}{11.1. Bộ dò đường Path Finder của Praat 7.0.03 và nguyên nhân loại bỏ Viterbi trên FPGA}

\subsubsection*{Vấn đề thực tế: Hiện tượng nhảy quãng tám cục bộ và sự phụ thuộc thời gian}
Giả sử hệ thống phần cứng hoàn tất việc tìm đỉnh tự tương quan, nội suy parabol và quyết định vô thanh / hữu thanh cho một khung âm thanh độc lập như đã thiết kế từ Chương 6 đến Chương 10. Tại mỗi khung phân tích cách nhau $10\text{ ms}$, mạch xuất ra một giá trị tần số cơ bản $F_{0,\text{Q12.4}}$ (hoặc gán bằng $0$ nếu khung là vô thanh). Câu hỏi đặt ra là: \emph{Nếu chỉ ghép nối độc lập các giá trị $F_0$ của từng khung rời rạc lại với nhau, đường biểu diễn cao độ theo thời gian trông sẽ như thế nào?}

Thực tế ghi âm tiếng nói cho thấy luồng âm thanh luôn bị tác động bởi nhiễu môi trường, hiện tượng cộng hưởng khoang miệng, hoặc sự biến thiên đột ngột của thanh môn (hiện tượng khàn tiếng, glottalization). Trong một khung âm thanh hữu thanh có giọng trầm $150\text{ Hz}$ ($\tau_0 \approx 107\text{ mẫu}$), hàm tự tương quan không chỉ xuất hiện một đỉnh nhọn tại $\tau_0$, mà còn có đỉnh thứ hai tại chu kỳ đôi $2\tau_0 \approx 214\text{ mẫu}$ (tương ứng với tần số chia đôi quãng tám $75\text{ Hz}$) và đỉnh tại $\tau_0 / 2 \approx 53\text{ mẫu}$ (nhân đôi quãng tám $300\text{ Hz}$).

Nếu khung âm thanh đó bị nhiễu trắng ngẫu nhiên chèn vào, chỉ cần một vài mẫu nhiễu làm biên độ đỉnh tương quan tại $2\tau_0$ nhỉnh hơn đỉnh thật $\tau_0$ một lượng siêu nhỏ (ví dụ chênh lệch $0{,}001$), bộ tìm đỉnh độc lập từng khung sẽ lập tức chọn đỉnh $75\text{ Hz}$. Hệ quả là trên một đoạn nguyên âm đang ngân dài êm ái ở mức $150\text{ Hz}$, đường cao độ đột ngột bị sụp xuống $75\text{ Hz}$ trong đúng 1 khung ($10\text{ ms}$), rồi nhịp sau lại nhảy ngược về $150\text{ Hz}$. Người nghe sẽ cảm nhận thấy một tiếng ``tách'' hoặc gián đoạn rất khó chịu. Trong âm học thực nghiệm, hiện tượng này được gọi là \textbf{lỗi nhảy quãng tám} (\emph{octave jump} hay \emph{pitch halving / pitch doubling}).

Để khắc phục hiện tượng nhảy bậc cục bộ này, phần mềm Praat không bao giờ xuất ngay kết quả của từng khung. Thay vào đó, Paul Boersma đã tích hợp một thuật toán quy hoạch động toàn cục có tên là \textbf{Bộ dò đường} (\emph{Path Finder}), dựa trên nguyên lý giải mã Viterbi nổi tiếng.

\subsubsection*{Thuật toán Viterbi Path Finder trong mã nguồn Praat 7.0.03}
Trong mã nguồn Praat (\path{fon/Sound_to_Pitch.cpp}, dòng 489--491),
hàm phân tích cao độ luôn gọi hàm \texttt{Pitch\_pathFinder}:
\begin{quote}
\small
\texttt{Pitch\_pathFinder (thee, silenceThreshold, voicingThreshold,}\\
\texttt{\ \ \ \ octaveCost, octaveJumpCost, voicedUnvoicedCost, ceiling);}
\end{quote}
Định nghĩa chi tiết của giải thuật này nằm tại tệp \texttt{fon/Pitch.cpp}, các dòng 524--652. Thay vì chỉ lưu một ứng viên tốt nhất, tại mỗi khung phân tích $t$ ($t = 1..N$), Praat giữ lại một danh sách tối đa gồm \textbf{15 ứng viên cao độ} (\texttt{maxnCandidates = 15}, dòng 548):
\begin{itemize}
    \item Ứng viên 1: Luôn là ứng viên vô thanh ($f = 0\text{ Hz}$, dòng 170--179 của \texttt{Sound\_to\_Pitch.cpp}).
    \item Ứng viên 2 đến 15: Là các đỉnh cực đại cục bộ của hàm tự tương quan đã được tinh chỉnh độ sâu bằng sinc70 và Brent.
\end{itemize}

Để tìm ra một đường đi xuyên suốt toàn bộ tệp âm thanh sao cho tổng chi phí là nhỏ nhất (hoặc tổng điểm số tích lũy là lớn nhất), Praat thiết lập đồ hình Trellis (Hình \ref{fig:ch11_viterbi_trellis}) và khởi tạo hai ma trận lớn tại các dòng 548--549 của \texttt{fon/Pitch.cpp}:
\begin{quote}
\small
\texttt{autoMAT delta = zero\_MAT (my nx, maxnCandidates);}\\
\texttt{autoINTMAT psi = zero\_INTMAT (my nx, maxnCandidates);}
\end{quote}
trong đó:
\begin{itemize}
    \item Ma trận \texttt{delta[t][i]} lưu trữ điểm số tích lũy tối ưu lớn nhất đi từ khung đầu tiên đến ứng viên thứ $i$ của khung thứ $t$.
    \item Ma trận \texttt{psi[t][i]} lưu con trỏ truy vết ngược (\emph{backpointer}), ghi nhận xem ứng viên $i$ ở khung $t$ được nối tới ứng viên nào ở khung $t-1$ để đạt điểm số cực đại đó.
\end{itemize}

\subsubsection*{Dẫn xuất các thành phần chi phí trong giải thuật Viterbi}
Tại mỗi nút của đồ hình, điểm số bao gồm hai thành phần: \textbf{chi phí cục bộ} (\emph{local cost}) và \textbf{chi phí chuyển tiếp} (\emph{transition cost}).

\paragraph{1. Chi phí cục bộ (Local cost).}
Tại \texttt{fon/Pitch.cpp} các dòng 553--561, điểm số cục bộ của từng ứng viên được thiết lập:
\begin{itemize}
    \item Đối với ứng viên 1 (vô thanh): nhận độ mạnh vô thanh $R_{\text{unvoiced}}$ (đã được phân tích chi tiết tại Chương 10):
    \begin{equation}
        \text{local\_cost}(c_1) = R_{\text{unvoiced}} = 0{,}45 + \max\left(0,\, 2 - \frac{LP}{GP} \cdot \frac{145}{3}\right)
    \end{equation}
    \item Đối với các ứng viên hữu thanh $c_i$ ($i = 2..15$, tần số $f_i$ nằm trong dải $[75, 600]\text{ Hz}$): nhận độ mạnh tương quan $s_i$ trừ đi chi phí quãng tám:
    \begin{equation}
        \text{local\_cost}(c_i) = s_i - \text{octaveCost} \cdot \log_2\left(\frac{\text{ceiling}}{f_i}\right)
    \end{equation}
    với $\text{octaveCost} = 0{,}01$ và $\text{ceiling} = 600\text{ Hz}$. Đại lượng này chính là điểm số $S_c$ chuẩn hóa mà mạch phần cứng của ta đã tính toán ở Chương 8 và Chương 9.
\end{itemize}

\paragraph{2. Chuẩn hóa bước thời gian (Time step correction).}
Tại \texttt{fon/Pitch.cpp} các dòng 543--545, các hệ số phạt chuyển tiếp được tỉ lệ hóa theo bước nhảy thời gian giữa hai khung liên tiếp:
\begin{equation}
    \text{timeStepCorrection} = \frac{0{,}01}{\Delta t}
\end{equation}
Với cấu hình chuẩn $f_s = 16.000\text{ Hz}$ và $\text{pitchFloor} = 75\text{ Hz}$, bước nhảy thời gian tự động của Praat là $\Delta t = 0{,}75 / 75 = 0{,}01\text{ s} = 10\text{ ms}$ (hop 160 mẫu, dòng 296 của \texttt{Sound\_to\_Pitch.cpp}). Khi đó:
\begin{equation}
    \text{timeStepCorrection} = \frac{0{,}01}{0{,}01} = 1{,}0
\end{equation}
Các hệ số chi phí chuyển tiếp giữ nguyên giá trị cấu hình danh định:
\begin{align}
    \text{currentOctaveJumpCost} &= 0{,}35 \times 1{,}0 = 0{,}35, \\
    \text{currentVoicedUnvoicedCost} &= 0{,}14 \times 1{,}0 = 0{,}14.
\end{align}

\paragraph{3. Chi phí chuyển trạng thái (Transition cost).}
Khi chuyển từ ứng viên $c_1$ (tần số $f_1$) ở khung $t-1$ sang ứng viên $c_2$ (tần số $f_2$) ở khung $t$, hàm phạt $\text{transitionCost}(c_1 \to c_2)$ tại các dòng 564--628 được tính như sau:
\begin{enumerate}
    \item \textbf{Chuyển tiếp Vô thanh $\to$ Vô thanh (dòng 586):} Cả hai khung cùng là khoảng lặng hoặc phụ âm xát. Không có sự thay đổi trạng thái dao động dây thanh, chi phí bằng $0$:
    \begin{equation}
        \text{transitionCost} = 0{,}0
    \end{equation}
    \item \textbf{Chuyển tiếp Đổi trạng thái Hữu thanh $\leftrightarrow$ Vô thanh (dòng 586, 590):} Khi tiếng nói chuyển từ nguyên âm sang phụ âm vô thanh hoặc ngược lại, sự chuyển pha này bị phạt một lượng cố định:
    \begin{equation}
        \text{transitionCost} = \text{currentVoicedUnvoicedCost} = 0{,}14
    \end{equation}
    Hàm phạt này ngăn chặn hiện tượng chập chờn bật/tắt cờ hữu thanh tại các khung chuyển tiếp.
    \item \textbf{Chuyển tiếp Hữu thanh $\to$ Hữu thanh (dòng 606):} Khi cả hai khung đều có tiếng nói, sự thay đổi tần số được đo bằng khoảng cách cao độ logarit (tính theo quãng tám). Hàm phạt tỉ lệ thuận với độ lệch quãng:
    \begin{equation}
        \text{transitionCost} = \text{currentOctaveJumpCost} \cdot \left| \log_2\left(\frac{f_1}{f_2}\right) \right| = 0{,}35 \cdot \left| \log_2\left(\frac{f_1}{f_2}\right) \right|
    \end{equation}
    Nếu tần số nhảy gấp đôi ($f_2 = 2f_1$), độ lệch là $|\log_2(2)| = 1{,}0$, chi phí chuyển trạng thái bị phạt nặng tới $0{,}35$ điểm!
\end{enumerate}

Lưu ý rằng trong tệp \texttt{fon/Pitch.cpp} tại các dòng 591--604 có một khối mã kiểm tra nhánh \texttt{if (Melder\_debug == 30)}. Đây chỉ là đoạn mã phục vụ thực nghiệm gỡ lỗi nội bộ của các tác giả Praat và mặc định luôn luôn tắt (\emph{off by default}) trong mọi bản phát hành chính thức.

\paragraph{4. Bước tính xuôi (Forward pass) và quy tắc phá vỡ thế hòa (Tie-breaking).}
Tại dòng 609, điểm số ứng viên mới được cập nhật theo nguyên lý Bellman:
\begin{equation}
    \text{value} = \delta[t-1][i_1] - \text{transitionCost}(i_1 \to i_2) + \text{local\_cost}[t][i_2]
\end{equation}
Praat tìm giá trị cực đại trên toàn bộ các ứng viên $i_1 \in [1, 15]$ của khung trước:
\begin{equation}
    \delta[t][i_2] = \max_{i_1} (\text{value}), \qquad \psi[t][i_2] = \arg\max_{i_1} (\text{value})
\end{equation}
Tại các dòng 613--623, phép so sánh lựa chọn cực đại sử dụng điều kiện \textbf{lớn hơn chặt} (\texttt{if (value > maximum)}). Điều này dẫn tới một quy tắc quan trọng: \emph{khi có hai đường đi cho điểm số bằng nhau tuyệt đối, ứng viên gặp trước (chỉ số $i_1$ nhỏ hơn) sẽ luôn được giữ lại}.

\paragraph{5. Kết thúc đường đi và Bước truy vết ngược (Backtracking).}
Khi tính toán tới khung cuối cùng $N$, Praat duyệt qua 15 ứng viên của khung $N$ để tìm điểm xuất phát có tổng điểm số tích lũy lớn nhất (dòng 630--639):
\begin{quote}
\small
\texttt{integer place = 1;}\\
\texttt{double maximum = delta [my nx] [1];}\\
\texttt{for (integer icand = 2; icand <= nCandidates; icand ++) \{}\\
\texttt{\ \ \ \ if (delta [my nx] [icand] > maximum) \{}\\
\texttt{\ \ \ \ \ \ \ \ place = icand; maximum = delta [my nx] [icand];}\\
\texttt{\ \ \ \ \}}\\
\texttt{\}}
\end{quote}
Biến \texttt{place} được khởi tạo bằng $1$ (ứng viên vô thanh) và chỉ cập nhật khi tìm được ứng viên khác có điểm số lớn hơn chặt. Do đó, nếu xảy ra hòa điểm ở khung cuối, \textbf{ứng viên vô thanh sẽ luôn chiến thắng}.

Sau khi xác định được \texttt{place} tại khung cuối $N$, thuật toán tiến hành \textbf{truy vết ngược} (\emph{backtracking}) từ khung $N$ ngược về khung 1 tại các dòng 641--652:
\begin{equation}
    \text{place}_{t-1} = \psi[t][\text{place}_t], \qquad \text{với } t = N, N-1, \dots, 2
\end{equation}
Tại mỗi khung, ứng viên tối ưu nằm trên đường đi toàn cục được hoán vị lên vị trí số 1 (\texttt{candidates[1]}), và tần số của nó được chọn làm kết quả $F_0$ chính thức của khung.

\begin{figure}[htbp]
    \centering
    % chuong_v3/fig_ch11_viterbi_trellis.tex
% Sơ đồ đồ hình Trellis của giải thuật Viterbi trong Praat Path Finder
% Đối chiếu giữa quy hoạch động toàn cục (Praat) và bộ lọc cục bộ Median-3 (FPGA)
\resizebox{\linewidth}{!}{%
\begin{tikzpicture}[
    >=Stealth,
    scale=0.88,
    frame_box/.style={rectangle, draw=blue!70!black, fill=blue!3, rounded corners=3pt, minimum width=1.6cm, minimum height=5.2cm, inner sep=2pt},
    cand_uv/.style={circle, draw=red!80!black, fill=red!10, thick, minimum size=0.68cm, inner sep=0pt, font=\scriptsize\bfseries},
    cand_v/.style={circle, draw=blue!80!black, fill=blue!10, thick, minimum size=0.68cm, inner sep=0pt, font=\scriptsize\bfseries},
    cand_best/.style={circle, draw=orange!90!black, fill=orange!25, very thick, minimum size=0.72cm, inner sep=0pt, font=\scriptsize\bfseries},
    wire_trans/.style={draw=black!30, thin, ->},
    wire_viterbi/.style={draw=red!85!black, very thick, ->},
    lbl_tag/.style={font=\scriptsize, fill=white, inner sep=2pt, rounded corners=2pt}
]

    % -------------------------------------------------------------------------
    % CÁC KHUNG THỜI GIAN TRONG ĐỒ HÌNH TRELLIS
    % -------------------------------------------------------------------------
    % Khung t-1
    \node[frame_box] (f_prev) at (2.0, 2.5) {};
    \node[below=7pt, font=\scriptsize\bfseries, text=blue!80!black] at (f_prev.south) {Khung $t-1$};

    % Khung t
    \node[frame_box] (f_cur) at (6.0, 2.5) {};
    \node[below=7pt, font=\scriptsize\bfseries, text=blue!80!black] at (f_cur.south) {Khung $t$};

    % Khung t+1
    \node[frame_box] (f_next) at (10.0, 2.5) {};
    \node[below=7pt, font=\scriptsize\bfseries, text=blue!80!black] at (f_next.south) {Khung $t+1$};

    % Khung N (Khung cuối cùng)
    \node[frame_box] (f_end) at (14.0, 2.5) {};
    \node[below=7pt, font=\scriptsize\bfseries, text=blue!80!black] at (f_end.south) {Khung cuối $N$};

    % Dấu chấm lửng giữa t+1 và N
    \node[font=\large\bfseries, text=black!60] at (12.0, 2.5) {$\cdots$};

    % -------------------------------------------------------------------------
    % CÁC ỨNG VIÊN TRONG TỪNG KHUNG (15 ứng viên mỗi khung)
    % -------------------------------------------------------------------------
    % Khung t-1
    \node[cand_uv] (c_p_uv) at (2.0, 4.4) {$c_1$};
    \node[cand_best] (c_p_v1) at (2.0, 3.1) {$c_2$};
    \node[cand_v] (c_p_v2) at (2.0, 1.8) {$c_3$};
    \node[font=\scriptsize, text=black!60] at (2.0, 0.7) {$\vdots$};
    \node[cand_v] (c_p_vm) at (2.0, 0.1) {$c_{15}$};

    % Khung t
    \node[cand_uv] (c_c_uv) at (6.0, 4.4) {$c_1$};
    \node[cand_best] (c_c_v1) at (6.0, 3.1) {$c_2$};
    \node[cand_v] (c_c_v2) at (6.0, 1.8) {$c_3$};
    \node[font=\scriptsize, text=black!60] at (6.0, 0.7) {$\vdots$};
    \node[cand_v] (c_c_vm) at (6.0, 0.1) {$c_{15}$};

    % Khung t+1
    \node[cand_uv] (c_n_uv) at (10.0, 4.4) {$c_1$};
    \node[cand_best] (c_n_v1) at (10.0, 3.1) {$c_2$};
    \node[cand_v] (c_n_v2) at (10.0, 1.8) {$c_3$};
    \node[font=\scriptsize, text=black!60] at (10.0, 0.7) {$\vdots$};
    \node[cand_v] (c_n_vm) at (10.0, 0.1) {$c_{15}$};

    % Khung N
    \node[cand_best] (c_e_uv) at (14.0, 4.4) {$c_1$};
    \node[cand_v] (c_e_v1) at (14.0, 3.1) {$c_2$};
    \node[cand_v] (c_e_v2) at (14.0, 1.8) {$c_3$};
    \node[font=\scriptsize, text=black!60] at (14.0, 0.7) {$\vdots$};
    \node[cand_v] (c_e_vm) at (14.0, 0.1) {$c_{15}$};

    % -------------------------------------------------------------------------
    % CÁC ĐƯỜNG CHUYỂN TRẠNG THÁI VÀ CHI PHÍ (TRANSITION COSTS)
    % -------------------------------------------------------------------------
    % 1. Chuyển tiếp vô thanh - vô thanh: UV -> UV (chi phí = 0)
    \draw[wire_trans, draw=green!60!black] (c_p_uv.east) -- (c_c_uv.west);
    \node[font=\scriptsize, fill=white, inner sep=2pt, rounded corners=2pt, draw=green!60!black, text=green!60!black] at (4.0, 4.95) {UV--UV: $0$};

    % 2. Đổi trạng thái V/UV: c_1 (vô thanh) -> c_2 (hữu thanh) (chi phí = 0,14)
    \draw[wire_trans, draw=purple!80!black] (c_p_uv.east) -- (c_c_v1.west);
    \path (c_p_uv.east) -- (c_c_v1.west) node[pos=0.48, sloped, above=3pt, font=\scriptsize, fill=white, inner sep=2pt, rounded corners=2pt, draw=purple!60, text=purple!80!black] {Đổi V/UV: $0{,}14$};

    % 3. Nhảy tần số F0: c_2 -> c_3 (chi phí = 0,35*|log2(f1/f2)|)
    \draw[wire_trans, draw=blue!80!black] (c_p_v1.east) -- (c_c_v2.west);
    \node[font=\scriptsize, fill=white, inner sep=2pt, rounded corners=2pt, draw=blue!50, text=blue!80!black, align=center] (lbl_nhay) at (4.0, 1.05) {Nhảy $F_0$:\\$0{,}35|\log_2\frac{f_1}{f_2}|$};
    \draw[draw=blue!60, thin, ->] (lbl_nhay.north) -- (4.0, 2.45);

    % Chuyển tiếp giữa khung t và t+1
    \draw[wire_trans] (c_c_uv.east) -- (c_n_uv.west);
    \draw[wire_trans] (c_c_v1.east) -- (c_n_v2.west);
    \draw[wire_trans] (c_c_v2.east) -- (c_n_v1.west);

    % Đường đi tối ưu Viterbi toàn cục (màu đỏ đậm có mũi tên)
    \draw[wire_viterbi] (c_p_v1.east) -- (c_c_v1.west);
    \draw[wire_viterbi] (c_c_v1.east) -- (c_n_v1.west);

    % -------------------------------------------------------------------------
    % ĐƯỜNG TRUY VẾT NGƯỢC (BACKTRACKING)
    % Mũi tên trỏ từ khung cuối N ngược về khung đầu (phải sang trái)
    % -------------------------------------------------------------------------
    \draw[draw=orange!90!black, ultra thick, dashed, ->] (14.0, 5.6) -- node[above=2pt, font=\scriptsize\bfseries, text=orange!90!black] {Truy vết ngược toàn bộ file ($N \to 1$)} (2.0, 5.6);

    % -------------------------------------------------------------------------
    % CHÚ THÍCH PHÂN TÁCH VITERBI VS MEDIAN-3
    % -------------------------------------------------------------------------
    \node[draw=red!70!black, fill=red!3, rounded corners=2pt, font=\scriptsize, align=left, inner sep=4pt] at (8.0, -1.6) {
        \textbf{Praat Path Finder (R1):} Lưu trữ toàn bộ ma trận $\delta, \psi$ kích thước $N \times 15$; bước tính xuôi cộng dồn chi phí chuyển tiếp,\\
        sau đó truy vết ngược từ khung cuối $N$ về khung 1 $\implies$ Bùng nổ RAM và độ trễ toàn file trên FPGA.
    };

\end{tikzpicture}%
}

    \caption{Đồ hình Trellis của giải thuật quy hoạch động Viterbi trong Praat Path Finder. Tại mỗi khung thời gian $t$, thuật toán duy trì 15 trạng thái ứng viên cạnh tranh. Đường đi tối ưu toàn cục (màu đỏ) được xác định bằng bước tính xuôi tích lũy chi phí và truy vết ngược từ khung cuối $N$ về khung 1.}
    \label{fig:ch11_viterbi_trellis}
\end{figure}

\textbf{Đọc hình.} Hình \ref{fig:ch11_viterbi_trellis} minh họa trực quan cấu trúc đồ hình Trellis của thuật toán Viterbi trong Praat. Mỗi cột chữ nhật biểu diễn một khung phân tích $10\text{ ms}$, chứa 15 nút tròn tương ứng với 15 ứng viên cao độ (nút đỏ $c_1$ là ứng viên vô thanh $0\text{ Hz}$, các nút xanh $c_2..c_{15}$ là các đỉnh tự tương quan hữu thanh). Các đường nét mờ thể hiện mạng lưới liên kết chuyển trạng thái giữa các khung kề nhau, kèm theo các hàm phạt nhảy tần số ($0{,}35|\log_2(f_1/f_2)|$) và đổi trạng thái vô thanh / hữu thanh ($0{,}14$). Đường mũi tên đỏ đậm là quỹ đạo cao độ toàn cục tối ưu được lựa chọn. Đường nét đứt màu cam ở phía trên thể hiện bước truy vết ngược (\emph{backtracking}): thuật toán bắt buộc phải đợi đến khi tính xong khung cuối cùng $N$ mới có thể lần ngược các con trỏ $\psi$ về khung 1 để chốt kết quả cao độ của từng khung.

\subsubsection*{Vì sao thiết kế phần cứng FPGA loại bỏ bộ dò đường Viterbi?}
Mặc dù Viterbi Path Finder đem lại đường pitch cực kỳ trơn tru trên phần mềm máy tính, việc hiện thực hóa nguyên bản thuật toán này trên chip FPGA Cyclone V vấp phải ba rào cản kiến trúc không thể vượt qua:

\begin{enumerate}
    \item \textbf{Bùng nổ tài nguyên bộ nhớ lưu trữ nội chip (On-chip RAM explosion):}
    Để thực hiện quy hoạch động Viterbi, FPGA bắt buộc phải lưu trữ đồng thời toàn bộ hai ma trận $\delta$ và $\psi$ trong bộ nhớ SRAM. Với một tệp âm thanh chỉ dài $10\text{ giây}$ ($1.000\text{ khung}$), số lượng phần tử cần lưu trữ là:
    \begin{equation}
        1.000\text{ khung} \times 15\text{ ứng viên} = 15.000\text{ trạng thái}
    \end{equation}
    Mỗi trạng thái bao gồm điểm số tích lũy (số thực hoặc số nguyên 32-bit) và con trỏ truy vết (4-bit). Ma trận dữ liệu đòi hỏi hơn $60\text{ KB}$ bộ nhớ. Nếu tệp âm thanh dài vài chục giây đến vài phút (ví dụ một câu chuyện dài 1 phút gồm $6.000\text{ khung}$), dung lượng RAM yêu cầu sẽ vượt quá $360\text{ KB}$ ($2{,}88\text{ Mbit}$).
    Trong khi đó, chip Cyclone V 5CSXFC6D6F31C6N chỉ có tổng cộng 553 khối M10K (tương đương khoảng $5{,}5\text{ Mbit}$ RAM), vốn phải chia sẻ cho các bộ đệm FIFO, RAM vòng lấy mẫu âm thanh, và các bảng sóng ROM Hanning, C/K. Việc dành trọn tài nguyên RAM nội chip chỉ để lưu ma trận Trellis cho một tệp có độ dài hữu hạn là phương án thiết kế không thực tế.

    \item \textbf{Phá vỡ hoàn toàn tính nhân quả và thông lượng thời gian thực (Non-causal latency):}
    Bản chất của bước truy vết ngược (\emph{backtracking}) là phải duyệt ngược chiều thời gian: từ khung $N$ lùi dần về khung 1. Điều này đồng nghĩa với việc phần cứng \textbf{không thể xuất ra bất kỳ giá trị $F_0$ nào cho đến khi mẫu âm thanh cuối cùng của tệp được nạp vào}!
    Độ trễ xử lý (latency) bị phụ thuộc vào độ dài của tệp âm thanh. Trong các ứng dụng đàm thoại trực tiếp, truyền âm thanh qua mạng hoặc hỗ trợ phát âm thời gian thực, độ trễ kéo dài hàng giây là hoàn toàn không thể chấp nhận được. Một hệ thống phần cứng chuyên dụng trên FPGA phải hoạt động theo cơ chế dòng (\emph{streaming}): mẫu âm thanh vào đến đâu, cao độ được tính toán và đẩy ra đến đó với độ trễ cố định, có thể dự đoán được ở mức chu kỳ xung nhịp.

    \item \textbf{Độ phức tạp điều khiển và quản lý ứng viên:}
    Việc duy trì 15 ứng viên đòi hỏi mạch phần cứng phải có bộ nhớ sắp xếp danh sách (sorting list), bộ chia và bộ tính logarit để lượng hóa chi phí chuyển tiếp $0{,}35|\log_2(f_1/f_2)|$ cho $15 \times 15 = 225$ cặp chuyển tiếp giữa mỗi cặp khung kề nhau. Lượng logic điều khiển và tài nguyên ALM tiêu hao cho khối Viterbi sẽ lấn át toàn bộ các khối xử lý tín hiệu số (DSP) cốt lõi của hệ thống.
\end{enumerate}

Cũng cần nhấn mạnh một sự thật kỹ thuật khách quan: \emph{trong chế độ xử lý tệp ngoại tuyến (file-based processing), bộ xử lý nhúng ARM Cortex-A9 của hệ thống HPS hoàn toàn có đủ năng lực bộ nhớ SDRAM để chạy giải thuật Path Finder bằng phần mềm nếu người dùng mong muốn}.
Tuy nhiên, mục tiêu của lõi xử lý FPGA trong tài liệu này là xây dựng một đường ống xử lý cao độ thuần phần cứng, độc lập, tối ưu hóa diện tích silicon và bảo đảm tính nhân quả thời gian thực.
Vì vậy, kiến trúc phần cứng của ta đã chốt quyết định: \textbf{loại bỏ hoàn toàn bộ dò đường Viterbi trên FPGA, thay thế bằng cơ chế chọn cực đại độc lập từng khung kết hợp với bộ lọc trung vị Median-3}.

Tài liệu này cũng kiên quyết loại bỏ khẳng định cảm tính vô căn cứ cho rằng ``Viterbi chỉ giúp ích dưới 3\% vì hơn 97\% số khung đã đúng ứng viên 1''. Con số 97\% là một suy đoán bịa đặt. Tầm quan trọng thực sự của Path Finder phụ thuộc hoàn toàn vào mức độ nhiễu của môi trường ghi âm; điều này sẽ được chứng minh bằng các con số thực nghiệm cụ thể và khách quan ở mục 11.2 ngay sau đây.

\textit{Tóm lại: Praat sử dụng giải thuật quy hoạch động Viterbi trên 15 ứng viên mỗi khung với hai ma trận $\delta, \psi$ để triệt tiêu lỗi nhảy quãng tám; tuy nhiên bước truy vết ngược phi nhân quả và nguy cơ cạn kiệt bộ nhớ RAM nội chip khiến Viterbi không phù hợp với kiến trúc luồng phần cứng trên FPGA.}

% ============================================================================
% 11.2. HAI CẤU HÌNH ĐỐI CHUẨN R1, R2 VÀ CHUẨN NGANG HÀNG
% ============================================================================
\subsection*{11.2. Hai cấu hình đối chuẩn R1, R2 và vai trò chuẩn ngang hàng của R2}
\addcontentsline{toc}{subsection}{11.2. Hai cấu hình đối chuẩn R1, R2 và vai trò chuẩn ngang hàng của R2}

\subsubsection*{Vấn đề thực tế: So sánh FPGA với Praat bằng chuẩn nào cho công bằng?}
Khi đánh giá độ chính xác của đường cao độ $F_0$ do chip FPGA tính toán so với phần mềm Praat, một câu hỏi phương pháp luận cốt lõi được đặt ra: \emph{Nếu ta so trực tiếp ngõ ra của FPGA với kết quả mặc định của Praat, phần sai lệch đo được đến từ đâu?}

Có phải phần chênh lệch đó đến từ sai số làm tròn số nguyên Q-format của FPGA? Hay đến từ phép nội suy parabol thay cho sinc70? Hay thực chất phần lớn độ lệch lại đến từ việc \emph{Praat có bộ dò đường Viterbi toàn cục còn FPGA thì không}?
Nếu đem một hệ thống không có Viterbi so sánh với một hệ thống có Viterbi đầy đủ, ta đang thực hiện một phép so sánh khập khiễng giữa ``táo và cam'' (\emph{apples to oranges}). Để bóc tách tường minh từng nguồn sai số, ta cần một phiên bản Praat tham chiếu đóng vai trò là \textbf{chuẩn đối sánh ngang hàng} (\emph{apples-to-apples}).

\subsubsection*{Sự suy biến của Viterbi khi chi phí chuyển trạng thái triệt tiêu về 0}
Hãy khảo sát điều gì sẽ xảy ra với giải thuật Viterbi của Praat khi ta can thiệp vào các tham số chi phí chuyển trạng thái bằng cách đặt:
\begin{equation}
    \text{octaveJumpCost} = 0{,}0, \qquad \text{voicedUnvoicedCost} = 0{,}0
\end{equation}
Khi cả hai hệ số này bằng $0$, theo công thức tính toán tại các dòng 586--606 của \texttt{fon/Pitch.cpp}:
\begin{itemize}
    \item Chi phí chuyển tiếp vô thanh -- vô thanh: $\text{transitionCost} = 0{,}0$.
    \item Chi phí chuyển tiếp đổi trạng thái vô thanh $\leftrightarrow$ hữu thanh: $\text{transitionCost} = \text{currentVoicedUnvoicedCost} = 0{,}0$.
    \item Chi phí chuyển tiếp hữu thanh -- hữu thanh: $\text{transitionCost} = 0{,}0 \cdot |\log_2(f_1/f_2)| = 0{,}0$.
\end{itemize}
Như vậy, \textbf{mọi chi phí chuyển tiếp giữa bất kỳ cặp trạng thái nào đều triệt tiêu về đúng $0$ tuyệt đối}:
\begin{equation}
    \text{transitionCost}(i_1 \to i_2) \equiv 0, \quad \forall i_1, i_2 \in [1, 15]
\end{equation}

Thay kết quả này vào phương trình cập nhật điểm số Bellman tại dòng 609:
\begin{equation}
    \text{value}(i_1, i_2) = \delta[t-1][i_1] - 0 + \text{local\_cost}[t][i_2] = \delta[t-1][i_1] + \text{local\_cost}[t][i_2]
\end{equation}
Để tìm điểm số tích lũy lớn nhất cho ứng viên $i_2$ tại khung $t$:
\begin{equation}
    \delta[t][i_2] = \max_{i_1} \left( \delta[t-1][i_1] + \text{local\_cost}[t][i_2] \right)
\end{equation}
Vì số hạng $\text{local\_cost}[t][i_2]$ hoàn toàn độc lập với chỉ số $i_1$, ta có thể đưa nó ra ngoài dấu cực đại:
\begin{equation}
    \delta[t][i_2] = \left( \max_{i_1} \delta[t-1][i_1] \right) + \text{local\_cost}[t][i_2]
    \label{eq:viterbi_degenerate}
\end{equation}

Biểu thức \eqref{eq:viterbi_degenerate} mang ý nghĩa toán học cực kỳ sâu sắc:
Đại lượng $M_{t-1} = \max_{i_1} \delta[t-1][i_1]$ là một hằng số cố định giống nhau đối với mọi ứng viên $i_2$ tại khung $t$. Hơn nữa, con trỏ truy vết ngược:
\begin{equation}
    \psi[t][i_2] = \arg\max_{i_1} \delta[t-1][i_1] = \text{constant}, \quad \forall i_2
\end{equation}
mọi ứng viên ở khung $t$ đều trỏ về cùng một trạng thái tối ưu duy nhất của khung $t-1$!

Tại khung cuối cùng $N$, việc tìm ứng viên có điểm số tích lũy lớn nhất:
\begin{equation}
    \text{place}_N = \arg\max_{i_2} \delta[N][i_2] = \arg\max_{i_2} \left( M_{N-1} + \text{local\_cost}[N][i_2] \right) = \arg\max_{i_2} \text{local\_cost}[N][i_2]
\end{equation}
chính là việc chọn ứng viên có chi phí cục bộ lớn nhất tại chính khung $N$!
Khi truy vết ngược về khung $N-1$:
\begin{equation}
    \text{place}_{N-1} = \psi[N][\text{place}_N] = \arg\max_{i_1} \delta[N-1][i_1] = \arg\max_{i_1} \text{local\_cost}[N-1][i_1]
\end{equation}
Bằng phương pháp quy nạp toán học, tại mọi khung thời gian $t \in [1, N]$, ứng viên được chọn trên đường đi Viterbi thoái hóa chính là:
\begin{equation}
    \mathbf{\text{place}_t \equiv \arg\max_{i \in [1, 15]} \text{local\_cost}[t][i]}
    \label{eq:argmax_equivalence}
\end{equation}

\textbf{Định lý suy biến:} \emph{Khi triệt tiêu chi phí nhảy quãng tám và chi phí đổi trạng thái vô thanh về $0$, giải thuật quy hoạch động Viterbi của Praat suy biến hoàn toàn thành phép chọn cực đại độc lập từng khung (per-frame argmax) của hàm chi phí cục bộ $\text{local\_cost}$}.
Trường hợp hòa điểm được phân xử đồng nhất: ứng viên có chỉ số nhỏ hơn (xuất hiện trước) sẽ chiến thắng do quy tắc so sánh lớn hơn chặt (\texttt{>}).

\subsubsection*{Kiểm chứng số học sự tương đương giữa R2 và Per-Frame Argmax}
Để kiểm chứng tính đúng đắn của định lý suy biến trên dữ liệu thực tế, script kiểm định \texttt{verify\_v3/ch11\_viterbi\_audit.py} đã nạp thuật toán tham chiếu \texttt{PraatRawAc} (với $\text{octave\_jump}=0{,}0$, $\text{vuv}=0{,}0$) và so sánh trực tiếp từng khung với phép tính $\text{argmax}(\text{delta}[t])$ trên toàn bộ 277 khung của hai tệp âm thanh chuẩn trong \texttt{golden\_v2}:
\begin{itemize}
    \item Tệp âm thanh sạch \texttt{synth\_glide\_vuv\_16k.wav} ($277\text{ khung}$):
    \begin{equation}
        \max_{1 \le t \le 277} |F_{0,\text{R2}}[t] - F_{0,\text{argmax}}[t]| = 0{,}0000000000000000\text{ Hz}
    \end{equation}
    Số khung khớp tuyệt đối là \textbf{277 / 277 khung (100,00\%)}.
    \item Tệp âm thanh có nhiễu $5\text{ dB}$ \texttt{synth\_glide\_vuv\_16k\_snr5.wav} ($277\text{ khung}$):
    \begin{equation}
        \max_{1 \le t \le 277} |F_{0,\text{R2}}[t] - F_{0,\text{argmax}}[t]| = 0{,}0000000000000000\text{ Hz}
    \end{equation}
    Số khung khớp tuyệt đối cũng đạt \textbf{277 / 277 khung (100,00\%)}.
\end{itemize}
Độ lệch bằng $0$ đến mức chính xác dấu phẩy động 64-bit khẳng định chắc chắn rằng: \textbf{Cấu hình Praat R2 chính là bản sao toán học của phép chọn cực đại đơn khung}!

\subsubsection*{Hai cấu hình đối chuẩn: Praat R1 và Praat R2}
Từ kết luận trên, tài liệu này định nghĩa rõ hai mốc tham chiếu chuẩn tắc của Praat 7.0.03:
\begin{itemize}
    \item \textbf{Cấu hình R1 (Praat mặc định):} $\text{octaveJumpCost} = 0{,}35$, $\text{voicedUnvoicedCost} = 0{,}14$. Đây là cấu hình đầy đủ của phần mềm Praat thương mại chạy trên máy tính, sử dụng trọn vẹn giải thuật Viterbi toàn cục.
    \item \textbf{Cấu hình R2 (Praat chi phí 0):} $\text{octaveJumpCost} = 0{,}0$, $\text{voicedUnvoicedCost} = 0{,}0$. Toàn bộ 8 tham số còn lại (tần số lấy mẫu, dải tần 75--600 Hz, cửa sổ Hanning 638 mẫu, ngưỡng im lặng 0,03, ngưỡng vô thanh 0,45, nội suy sinc70 + Brent) được giữ nguyên vẹn như R1.
\end{itemize}

Ý nghĩa phương pháp luận của cặp đối chuẩn này trong hệ thống:
\begin{enumerate}
    \item \textbf{R2 là chuẩn ngang hàng (\emph{apples-to-apples}) cho ngõ ra FPGA thô (trước Median-3):}
    Cả hai bên đều thực hiện cùng một thuật toán: chọn ứng viên tốt nhất trên từng khung độc lập. Do đó, bất kỳ sự chênh lệch nào giữa FPGA và R2 chỉ có thể xuất phát từ các yếu tố phần cứng: phép nội suy parabol thay cho sinc70, số nguyên Q-format thay cho float64, và việc xếp hạng ứng viên trên lưới rời rạc thay cho đỉnh liên tục.
    \item \textbf{Khoảng lệch R1 -- R2 đo lường chính xác phần đóng góp của Path Finder:}
    Vì R1 và R2 sử dụng chung $100\%$ các bước xử lý tín hiệu phía trước (cùng cửa sổ, cùng tự tương quan, cùng danh sách 15 ứng viên), sự khác biệt duy nhất giữa R1 và R2 phản ánh chính xác hiệu quả mà giải thuật Viterbi mang lại cho bài toán.
\end{enumerate}

\subsubsection*{Phân tích định lượng đóng góp của Path Finder trên dữ liệu thực nghiệm}
Bảng \ref{tab:ch11_r1_vs_r2} tổng hợp kết quả đối chiếu giữa R1, R2 và mô hình phần cứng FPGA trên hai tệp âm thanh chuẩn của dự án (trích xuất từ nhật ký \texttt{verify\_v3/ch11\_viterbi\_audit.log}):

\begin{table}[htbp]
\centering
\small
\setlength{\tabcolsep}{5pt}
\caption{Bảng đối chiếu định lượng giữa Praat R1, Praat R2 và lõi phần cứng FPGA trên hai tệp kiểm chứng tổng hợp (277 khung, $f_s = 16\text{ kHz}$).}
\label{tab:ch11_r1_vs_r2}
\begin{tabular}{p{7.2cm}rr}
\toprule
\textbf{Chỉ số đo lường} & \textbf{Tệp sạch} & \textbf{Tệp nhiễu $5\text{ dB}$} \\
& \textbf{(\texttt{synth\_glide})} & \textbf{(\texttt{synth\_glide\_snr5})} \\
\midrule
Số khung phân tích toàn tệp & $277\text{ khung}$ & $277\text{ khung}$ \\
Số khung hữu thanh thực tế (Truth) & $207\text{ khung}$ & $207\text{ khung}$ \\
Lỗi quyết định vô thanh VDE giữa R1 và R2 & $0 / 277$ ($0{,}00\%$) & $0 / 277$ ($0{,}00\%$) \\
Số khung khác biệt tần số giữa R1 và R2 & \textbf{0 / 208 khung (0,00\%)} & \textbf{19 / 207 khung (9,18\%)} \\
Lỗi cao độ thô GPE của R1 so với F0 thật & \textbf{0 / 207 khung (0,00\%)} & \textbf{0 / 207 khung (0,00\%)} \\
Lỗi cao độ thô GPE của R2 so với F0 thật & \textbf{0 / 207 khung (0,00\%)} & \textbf{19 / 207 khung (9,18\%)} \\
Lỗi GPE của FPGA thô (trước Median) so với Truth & $0 / 207$ ($0{,}00\%$) & \textbf{22 / 207 khung (10,63\%)} \\
Lỗi GPE của FPGA sau Median-3 so với Truth & $0 / 207$ ($0{,}00\%$) & \textbf{15 / 207 khung (7,25\%)} \\
\bottomrule
\end{tabular}
\end{table}

\textbf{Đọc bảng.} Bảng \ref{tab:ch11_r1_vs_r2} cung cấp những hiểu biết vật lý vô cùng giá trị:
\begin{itemize}
    \item \textbf{Trên tín hiệu sạch (\texttt{synth\_glide\_vuv\_16k.wav}):}
    Các đỉnh tự tương quan của tín hiệu tuần hoàn thuần túy nổi bật rất rõ ràng trên nền nhiễu bằng 0. Điểm số cục bộ của ứng viên thật luôn vượt trội hoàn toàn so với các hài phụ. Kết quả là \textbf{R1 trùng khớp R2 ở $100\%$ số khung} ($0 / 277$ khung khác biệt). Trên tín hiệu sạch, bộ dò đường Viterbi hoàn toàn không phải can thiệp sửa chữa bất kỳ khung nào; việc chọn cực đại độc lập từng khung cho kết quả chính xác tuyệt đối.
    \item \textbf{Trên tín hiệu có nhiễu $5\text{ dB}$ (\texttt{synth\_glide\_vuv\_16k\_snr5.wav}):}
    Khi mức nhiễu tăng cao, năng lượng nhiễu ngẫu nhiên làm biến dạng cục bộ một số đỉnh tự tương quan. R1 khác R2 ở đúng \textbf{19 trên 207 khung hữu thanh} (tỉ lệ $9{,}18\%$). Khi đối chiếu với tần số thực tế của tín hiệu nguồn ($F_{0,\text{truth}}$):
    \begin{itemize}
        \item Cả 19 khung khác biệt này đều là \textbf{lỗi nhảy quãng tám của R2} (tần số bị sụp xuống một nửa hoặc một phần ba do đỉnh tự tương quan tại $2\tau$ hoặc $3\tau$ bị nhiễu đẩy cao hơn đỉnh chính). Do đó R2 có $\text{GPE} = 19 / 207$.
        \item Ngược lại, nhờ hàm phạt $\text{octaveJumpCost} = 0{,}35$, Praat R1 đã trừng phạt nặng nề các bước nhảy này và ép đường đi bám chặt vào nhánh tần số đúng, đạt kết quả hoàn hảo \textbf{$\text{GPE} = 0 / 207$ ($0{,}00\%$)}.
    \end{itemize}
    \item \textbf{Hiệu quả của lõi FPGA và Median-3:}
    Mô hình FPGA thô trước bộ lọc trung vị thừa hưởng các lỗi của phép chọn cực đại đơn khung, cộng thêm một số khung mấp mé ngưỡng lượng tử hóa nên đạt $\text{GPE} = 22 / 207$.
    Khi đưa qua bộ lọc trung vị Median-3, số khung lỗi GPE giảm từ \textbf{22 xuống 15 khung} (loại bỏ được 7 khung lỗi gai nhọn đơn lẻ). Tuy nhiên, Median-3 không thể triệt tiêu 15 khung lỗi còn lại do các lỗi này kéo dài từ 2 khung liên tiếp trở lên.
\end{itemize}

Con số giảm GPE từ $22$ xuống $15$ khung là minh chứng thực nghiệm không thể chối cãi cho kết luận: \emph{Median-3 là một bộ lọc làm mượt cục bộ hữu ích, nhưng nó không thể thay thế hoàn toàn năng lực sửa lỗi toàn cục của thuật toán Viterbi Path Finder}.

\textit{Tóm lại: Khi chi phí chuyển trạng thái bằng 0, Viterbi suy biến thành phép chọn cực đại đơn khung; do đó Praat R2 là chuẩn đối sánh ngang hàng trực tiếp cho phần cứng trước median, còn khoảng lệch R1 -- R2 đo lường chính xác 19 khung lỗi quãng tám mà Path Finder khắc phục được trên tệp nhiễu $5\text{ dB}$.}

% ============================================================================
% 11.3. BẢN CHẤT TOÁN HỌC CỦA BỘ LỌC MEDIAN-3
% ============================================================================
\subsection*{11.3. Bản chất toán học của bộ lọc Median-3: Những gì sửa được, không sửa được và tác động tới VDE}
\addcontentsline{toc}{subsection}{11.3. Bản chất toán học của bộ lọc Median-3: Những gì sửa được, không sửa được và tác động tới VDE}

\subsubsection*{Vấn đề thực tế: Vì sao chọn bộ lọc trung vị thay cho bộ lọc trung bình tuyến tính?}
Khi tín hiệu cao độ xuất hiện các xung lỗi đơn lẻ (ví dụ chuỗi $F_0$ qua các khung là $150\text{ Hz} \to 300\text{ Hz} \to 150\text{ Hz}$ do một khung bị nhân đôi quãng tám), giải pháp làm mượt thông thường trong xử lý tín hiệu là áp dụng một bộ lọc thông thấp số, chẳng hạn bộ lọc trung bình trượt 3 điểm (\emph{Moving Average}):
\begin{equation}
    y[k] = \frac{x[k-1] + x[k] + x[k+1]}{3}
\end{equation}
Nếu áp dụng bộ lọc trung bình cho chuỗi trên:
\begin{equation}
    y[1] = \frac{150 + 300 + 150}{3} = \frac{600}{3} = 200\text{ Hz}
\end{equation}
Kết quả thu được là giá trị $200\text{ Hz}$! Lỗi xung nhọn $300\text{ Hz}$ không hề biến mất mà bị ``nhoè'' sang các khung lân cận, làm sai lệch cả những khung vốn đang có giá trị đúng $150\text{ Hz}$.

Khác biệt căn bản của \textbf{bộ lọc trung vị} (\emph{Median Filter}) là tính chất phi tuyến: nó chọn giá trị nằm ở chính giữa sau khi đã sắp xếp thứ tự các phần tử trong cửa sổ quan sát. Đối với cửa sổ 3 điểm:
\begin{equation}
    y[k] = \text{median}(x[k-1],\, x[k],\, x[k+1])
\end{equation}
Sắp xếp ba giá trị $\{150, 300, 150\}$ thành dãy tăng dần: $\{150, 150, 300\}$. Phần tử trung vị ở giữa là đúng $150\text{ Hz}$! Xung lỗi $300\text{ Hz}$ bị triệt tiêu hoàn toàn mà không để lại bất kỳ vết nhòe năng lượng nào lên tín hiệu xung quanh.

\subsubsection*{Mạng logic so sánh 3 phần tử (Sorting Network)}
Trong toán học mạch số, phép tìm trung vị của 3 số nguyên không dấu 16-bit $a, b, c$ không đòi hỏi phải chạy các thuật toán sắp xếp phức tạp như Quicksort hay Bubble sort. Ta chỉ cần đúng \textbf{3 bộ so sánh nhị phân}:
\begin{equation}
    c_0 = (a \le b), \qquad c_1 = (b \le c), \qquad c_2 = (a \le c)
\end{equation}
Dựa trên kết quả logic của 3 bộ so sánh, giá trị trung vị được xác định thông qua cây ghép kênh MUX 3 nhánh:
\begin{equation}
    \text{median}(a, b, c) =
    \begin{cases}
        b & \text{nếu } (a \le b) \land (b \le c) \quad (a \le b \le c) \\
        c & \text{nếu } (a \le b) \land (c < b) \land (a \le c) \quad (a \le c < b) \\
        a & \text{nếu } (a \le b) \land (c < a) \quad (c < a \le b) \\
        a & \text{nếu } (b < a) \land (a \le c) \quad (b < a \le c) \\
        c & \text{nếu } (b < a) \land (c < a) \land (b \le c) \quad (b \le c < a) \\
        b & \text{nếu } (b < a) \land (c < b) \quad (c < b < a)
    \end{cases}
    \label{eq:median3_cases}
\end{equation}
Bằng cách gom nhóm điều kiện:
\begin{itemize}
    \item Nếu $a \le b$: nếu $b \le c$ thì lấy $b$; ngược lại nếu $a \le c$ thì lấy $c$; còn lại lấy $a$.
    \item Nếu $a > b$: nếu $a \le c$ thì lấy $a$; ngược lại nếu $b \le c$ thì lấy $c$; còn lại lấy $b$.
\end{itemize}
Mạng so sánh này bao phủ trọn vẹn $3! = 6$ hoán vị không gian cùng toàn bộ các trường hợp bằng nhau ($a = b$, $b = c$, hoặc $a = b = c$). Trễ lan truyền chỉ gồm 1 tầng so sánh trừ 16-bit và 1 tầng MUX tổ hợp, hoàn toàn đáp ứng thời gian xác lập trong chu kỳ xung nhịp $20\text{ ns}$ ($50\text{ MHz}$) của FPGA.

\subsubsection*{Kiểm chứng ví dụ tính tay thu nhỏ từ Đề cương}
Hãy đi qua từng bước tính toán thủ công trên dãy 6 khung F0 mẫu đã được chốt trong đề cương:
\begin{equation}
    \mathbf{x = [120,\, 122,\, 244,\, 121,\, 0,\, 0]}
\end{equation}
Theo quy tắc xử lý biên chuẩn của Praat và \texttt{golden\_v2}, dãy được đệm thêm mẫu lặp ở hai đầu biên:
\begin{equation}
    x_{\text{pad}} = [x_0,\, x_0,\, x_1,\, x_2,\, x_3,\, x_4,\, x_5,\, x_5] = [120,\, \mathbf{120,\, 122,\, 244,\, 121,\, 0,\, 0},\, 0]
\end{equation}
Áp dụng cửa sổ trung vị 3 điểm trượt qua từng vị trí $k = 0..5$:
\begin{itemize}
    \item Khung $k = 0$: Cửa sổ gồm $[120, 120, 122]$. Dãy sắp xếp $\{120, 120, 122\} \implies y[0] = \mathbf{120}$.
    \item Khung $k = 1$: Cửa sổ gồm $[120, 122, 244]$. Dãy sắp xếp $\{120, 122, 244\} \implies y[1] = \mathbf{122}$.
    \item Khung $k = 2$: Cửa sổ gồm $[122, 244, 121]$. Dãy sắp xếp $\{121, 122, 244\} \implies y[2] = \mathbf{122}$.
    \item Khung $k = 3$: Cửa sổ gồm $[244, 121, 0]$. Dãy sắp xếp $\{0, 121, 244\} \implies y[3] = \mathbf{121}$.
    \item Khung $k = 4$: Cửa sổ gồm $[121, 0, 0]$. Dãy sắp xếp $\{0, 0, 121\} \implies y[4] = \mathbf{0}$.
    \item Khung $k = 5$: Cửa sổ gồm $[0, 0, 0]$. Dãy sắp xếp $\{0, 0, 0\} \implies y[5] = \mathbf{0}$.
\end{itemize}
Dãy kết quả ngõ ra thu được là:
\begin{equation}
    \mathbf{y = [120,\, 122,\, 122,\, 121,\, 0,\, 0]}
\end{equation}
Tại khung $k = 2$, giá trị đột biến $244\text{ Hz}$ (lỗi nhân đôi quãng tám so với $122\text{ Hz}$) đã bị loại bỏ hoàn toàn, thay thế bằng giá trị mượt mà $122\text{ Hz}$. Kết quả tính tay này trùng khớp $100\%$ với nhật ký kiểm chứng \texttt{verify\_v3/ch11\_median\_math.log}.

\subsubsection*{Những gì Median-3 sửa được và không sửa được}
Việc nắm vững ranh giới năng lực của bộ lọc trung vị là yêu cầu bắt buộc đối với người thiết kế hệ thống nhúng:

\paragraph{1. Sửa được: Gai đơn lẻ một khung (Single spike).}
Khi một lỗi nhảy quãng tám chỉ xuất hiện cục bộ trong đúng 1 khung phân tích duy nhất và được bao bọc bởi hai khung có giá trị đúng ở hai bên:
\begin{equation}
    [\dots, F_0, \mathbf{2F_0}, F_0, \dots] \implies \text{median}(F_0, 2F_0, F_0) = F_0
\end{equation}
Phần tử lỗi $2F_0$ luôn bị đẩy về biên cực đại của tập 3 phần tử $\{F_0, F_0, 2F_0\}$, do đó giá trị trung vị luôn là $F_0$. Gai đơn bị triệt tiêu $100\%$.

\paragraph{2. Không sửa được: Gai kép kéo dài từ hai khung liên tiếp (Double spike).}
Giả sử nhiễu môi trường khiến lỗi nhân đôi quãng tám kéo dài liên tục trong 2 khung kề nhau:
\begin{equation}
    [\dots, 150, \mathbf{300, 300}, 150, \dots]
\end{equation}
Xét hai khung bị lỗi:
\begin{itemize}
    \item Tại khung lỗi thứ nhất: cửa sổ gồm $[150, \mathbf{300, 300}]$. Dãy sắp xếp $\{150, 300, 300\} \implies \text{trung vị} = \mathbf{300}$.
    \item Tại khung lỗi thứ hai: cửa sổ gồm $[\mathbf{300, 300}, 150]$. Dãy sắp xếp $\{150, 300, 300\} \implies \text{trung vị} = \mathbf{300}$.
\end{itemize}
Dãy kết quả ngõ ra vẫn giữ nguyên là $[150, \mathbf{300, 300}, 150]$!
Vì giá trị lỗi $300\text{ Hz}$ chiếm đa số ($2 / 3$ phần tử trong cửa sổ), bộ lọc trung vị 3 điểm hoàn toàn bất lực. Đây chính là lý do vì sao trên tệp âm thanh có nhiễu $5\text{ dB}$, Median-3 chỉ giảm được lỗi GPE từ 22 xuống 15 khung; 15 khung lỗi còn lại tồn tại dưới dạng các cặp gai kép mà chỉ có thuật toán tối ưu toàn cục như Viterbi mới có thể bẻ gãy.

\subsubsection*{Tác động của giá trị 0 lên quyết định vô thanh (VDE)}
Một đặc tính thiết kế quan trọng trong luồng xử lý của ta là: \textbf{giá trị $0$ biểu diễn trạng thái vô thanh tham gia trực tiếp vào phép so sánh trung vị}. Điều này mang lại cả mặt tích cực lẫn những tác động phụ lên chỉ số lỗi quyết định vô thanh (\emph{Voicing Decision Error} -- VDE):

\begin{enumerate}
    \item \textbf{Triệt tiêu đốm sáng giả trong khoảng lặng (Glitch suppression):}
    Nếu một khung im lặng ngẫu nhiên bị nhiễu kích hoạt nhầm cờ hữu thanh ($F_0 = 150\text{ Hz}$) giữa một chuỗi các khung vô thanh:
    \begin{equation}
        [0, 0, \mathbf{150}, 0, 0] \implies \text{median}(0, 150, 0) = \mathbf{0}
    \end{equation}
    Hai số 0 ở hai bên chiếm đa số, kéo giá trị $150\text{ Hz}$ trở về $0$. Đốm sáng giả bị xóa bỏ, giúp làm sạch khoảng lặng và cải thiện chất lượng nhận dạng tiếng nói.
    \item \textbf{Hàn gắn hiện tượng mất tiếng đơn khung (Dropout repair):}
    Ngược lại, nếu trong một đoạn nguyên âm hữu thanh ngân dài ổn định, có một khung bị sụt biên độ khiến bộ voicing ở Chương 10 ngắt nhầm về vô thanh ($F_0 = 0$):
    \begin{equation}
        [150, 150, \mathbf{0}, 150, 150] \implies \text{median}(150, 0, 150) = \mathbf{150}
    \end{equation}
    Hai giá trị $150\text{ Hz}$ ở hai bên chiếm đa số, hàn gắn khung mất tiếng trở lại thành $150\text{ Hz}$, khôi phục tính liên tục của nguyên âm.
    \item \textbf{Dịch chuyển nhẹ biên chuyển tiếp:}
    Tại điểm bắt đầu của một từ tiếng nói ($[0, 0, 150, 152, 154]$), tại khung hữu thanh đầu tiên cửa sổ là $[0, 150, 152] \implies \text{med} = 150$. Tuy nhiên tại khung trước đó $[0, 0, 150] \implies \text{med} = 0$. Median-3 có xu hướng làm sắc nét hoặc dịch chuyển biên vô thanh/hữu thanh khoảng 1 khung ($10\text{ ms}$).
\end{enumerate}

\textit{Tóm lại: Median-3 sử dụng mạng 3 bộ so sánh nhị phân để lọc phi tuyến; nó xóa sạch các gai lỗi đơn lẻ và hàn gắn các khung mất tiếng 0 cục bộ, nhưng hoàn toàn bất lực trước các lỗi kéo dài 2 khung liên tiếp.}

% ============================================================================
% 11.4. CĂN KHUNG NHÂN QUẢ PHẦN CỨNG
% ============================================================================
\subsection*{11.4. Căn khung nhân quả phần cứng, xử lý biên lặp mẫu và nhịp xả (Flush beat)}
\addcontentsline{toc}{subsection}{11.4. Căn khung nhân quả phần cứng, xử lý biên lặp mẫu và nhịp xả (Flush beat)}

\subsubsection*{Vấn đề thực tế: Tính phi nhân quả của bộ lọc trung vị căn giữa}
Trong công thức toán học và mô hình phần mềm tham chiếu (\texttt{fixed\_twin.py}, \texttt{praat\_float\_ref.py}), bộ lọc trung vị được định nghĩa ở dạng \textbf{căn giữa đối xứng} (\emph{centred}):
\begin{equation}
    y[k] = \text{median}(x[k-1],\, x[k],\, x[k+1])
\end{equation}
Để tính được giá trị $y[k]$ tại khung $k$, ta bắt buộc phải biết giá trị $x[k+1]$ của khung \emph{tương lai}.

Trên một hệ thống máy tính xử lý tệp, điều này rất đơn giản vì toàn bộ mảng dữ liệu đã nằm sẵn trong bộ nhớ. Nhưng trên mạch phần cứng thời gian thực FPGA, tại thời điểm đồng hồ đang xử lý khung $k$, khung $k+1$ \textbf{chưa hề được thu thập hay gửi tới}!
Nếu muốn tính trung vị mà không vi phạm nguyên lý nhân quả của vật lý, mạch phần cứng bắt buộc phải chấp nhận một sự đánh đổi: \textbf{độ trễ đường ống (latency) bằng đúng 1 khung}.

\subsubsection*{Căn khung nhân quả: Xuất khung $k-1$ tại thời điểm nhận khung $k$}
Giải pháp nhân quả chuẩn tắc trên phần cứng là làm trễ dòng dữ liệu 1 khung thời gian:
\begin{equation}
    y[k-1] = \text{median}(x[k-2],\, x[k-1],\, x[k])
\end{equation}
Trong chu kỳ thời gian khi khung mới $x[k]$ đang được đưa vào ngõ vào $\texttt{in\_f0}$:
\begin{itemize}
    \item Tín hiệu ngõ vào $\texttt{in\_f0}$: Mang giá trị mẫu mới nhất của khung hiện tại $x[k]$.
    \item Thanh ghi $\texttt{f0\_z1}$: Đang lưu giữ mẫu của khung vừa qua $x[k-1]$ (chính là khung trung tâm cần lọc).
    \item Thanh ghi $\texttt{f0\_z2}$: Đang lưu giữ mẫu quá khứ $x[k-2]$.
\end{itemize}
Mạng so sánh tổ hợp tính toán trung vị của bộ ba $\{\texttt{f0\_z2}, \texttt{f0\_z1}, \texttt{in\_f0}\}$ và chốt vào thanh ghi ngõ ra $\texttt{f0\_med}$ tại cạnh lên xung nhịp (\texttt{posedge clk}). Đồng thời, các thanh ghi cập nhật để chuẩn bị cho chu kỳ kế tiếp: $\texttt{f0\_z2} \leftarrow \texttt{f0\_z1}$ ($x[k-1]$), $\texttt{f0\_z1} \leftarrow \texttt{in\_f0}$ ($x[k]$), và $\texttt{f0\_z0} \leftarrow \texttt{in\_f0}$ ($x[k]$).
Do đó, trạng thái thanh ghi quan sát được sau mỗi cạnh lên xung nhịp (sau \texttt{posedge clk}) là: $\texttt{f0\_z0}$ và $\texttt{f0\_z1}$ cùng lưu mẫu mới nạp $x[k]$, $\texttt{f0\_z2}$ lưu mẫu $x[k-1]$, còn $\texttt{f0\_med}$ lưu kết quả trung vị đã lọc của khung $k-1$. Kết quả trung vị của khung $k-1$ được xuất ra tại thời điểm nhận khung $k$, với độ trễ nhân quả chính xác là 1 khung ($10\text{ ms}$).

\subsubsection*{Quy tắc xử lý biên lặp mẫu và nhịp xả tự động (Flush beat)}
Để kết quả nhân quả của phần cứng trùng khớp từng bit với công thức căn giữa của Praat trên toàn bộ $N$ khung, hai tình huống biên cần được xử lý đặc biệt:

\paragraph{1. Biên bắt đầu (Khung 0): Cơ chế lặp mẫu biên.}
Tại khung đầu tiên $k = 0$, tín hiệu điều khiển $\texttt{in\_first} = 1$ được kích hoạt cùng với $\texttt{in\_valid} = 1$. Vì chưa có các khung quá khứ $x[-1]$ hay $x[-2]$, phần cứng thực hiện quy tắc \textbf{lặp mẫu biên} (\emph{edge replication}) bằng cách nạp đồng thời giá trị $x[0]$ vào cả ba thanh ghi:
\begin{equation}
    \texttt{f0\_z2} \leftarrow x[0], \qquad \texttt{f0\_z1} \leftarrow x[0], \qquad \texttt{f0\_z0} \leftarrow x[0]
\end{equation}
Tại chu kỳ này, mạch hạ cờ ngõ ra $\texttt{out\_valid} = 0$ vì chưa có mẫu tương lai $x[1]$.
Khi khung tiếp theo $x[1]$ đến ở chu kỳ sau: mẫu $x[1]$ đóng vai trò là $x[k+1]$, mạch tính trung vị của $\{\texttt{f0\_z2}=x[0], \texttt{f0\_z1}=x[0], \texttt{in\_f0}=x[1]\}$ và xuất ra kết quả của khung 0:
\begin{equation}
    y[0] = \text{median}(x[0],\, x[0],\, x[1])
\end{equation}
hoàn toàn trùng khớp với công thức đệm biên phần mềm!

\paragraph{2. Biên kết thúc (Khung $N-1$): Nhịp xả tự động (Flush beat).}
Khi khung cuối cùng của tệp $x[N-1]$ được đưa vào mạch, tín hiệu $\texttt{in\_last} = 1$ bật lên. Tại chu kỳ này, mạch xuất ra kết quả của khung áp chót:
\begin{equation}
    y[N-2] = \text{median}(x[N-3],\, x[N-2],\, x[N-1])
\end{equation}
Lúc này, khung cuối cùng $x[N-1]$ được dịch vào thanh ghi trung tâm $\texttt{f0\_z1}$, còn khung $x[N-2]$ nằm ở $\texttt{f0\_z2}$. Nhưng vì tệp âm thanh đã hết, không còn khung $x[N]$ nào được gửi tới nữa! Kết quả của khung cuối $y[N-1]$ vẫn đang bị ``kẹt'' lại bên trong phần cứng.

Để giải phóng kết quả này, khối điều khiển FSM của module tự động kích hoạt một \textbf{nhịp xả} (\emph{Flush beat}) ngay ở chu kỳ xung nhịp kế tiếp. Trong nhịp xả:
\begin{itemize}
    \item Mạch thực hiện lặp mẫu cuối bằng cách đấu nối $x[N-1]$ (đang có sẵn trong $\texttt{f0\_z1}$) vào nhánh ngõ vào thứ ba của bộ so sánh.
    \item Mạng so sánh tính toán trung vị của bộ ba $\{x[N-2], x[N-1], x[N-1]\}$:
    \begin{equation}
        y[N-1] = \text{median}(x[N-2],\, x[N-1],\, x[N-1])
    \end{equation}
    \item Mạch kích hoạt đồng thời $\texttt{out\_valid} = 1$ và $\texttt{out\_last} = 1$ để thông báo đây là khung kết thúc của tệp.
\end{itemize}
Nhờ có nhịp xả, một tệp âm thanh có $N$ khung ngõ vào sẽ phát ra chính xác $N$ khung ngõ ra hoàn chỉnh, bảo đảm tính toàn vẹn dữ liệu tuyệt đối.

\begin{figure}[htbp]
    \centering
    % chuong_v3/fig_ch11_median_timing.tex
% Giản đồ định thời nhân quả của module ch11_median3
% Minh họa: Khung k-1 xuất tại chu kỳ k, lặp mẫu biên và nhịp xả Flush Beat
\resizebox{\linewidth}{!}{%
\begin{tikzpicture}[
    >=Stealth,
    scale=0.88,
    sig_name/.style={font=\scriptsize\ttfamily, anchor=east},
    grid_line/.style={draw=black!15, dashed},
    clk_line/.style={draw=black!85, thick},
    data_bus/.style={draw=blue!80!black, thick},
    val_line/.style={draw=green!60!black, thick},
    annot/.style={font=\scriptsize, fill=white, inner sep=1.5pt}
]

    % -------------------------------------------------------------------------
    % LƯỚI CHU KỲ XUNG NHỊP T1 .. T8
    % -------------------------------------------------------------------------
    \foreach \t in {1,2,3,4,5,6,7,8} {
        \draw[grid_line] (\t*2.0 - 0.5, 5.5) -- (\t*2.0 - 0.5, -5.3);
        \node[font=\scriptsize\bfseries, text=black!70] at (\t*2.0 + 0.5, 5.3) {T\t};
    }

    % Xung nhịp CLK (50 MHz)
    \node[sig_name] at (1.2, 4.6) {clk};
    \draw[clk_line] (1.5, 4.3) \foreach \i in {1,...,7} {
        -- ++(1.0, 0) -- ++(0, 0.6) -- ++(1.0, 0) -- ++(0, -0.6)
    };

    % -------------------------------------------------------------------------
    % CÁC TÍN HIỆU ĐẦU VÀO
    % -------------------------------------------------------------------------
    % in_valid
    \node[sig_name] at (1.2, 3.5) {in\_valid};
    \draw[val_line] (1.5, 3.0) -- (1.5, 3.4) -- (13.5, 3.4) -- (13.5, 3.0) -- (15.5, 3.0);
    \node[annot, text=green!50!black] at (7.5, 3.4) {Mức 1 (Khung 0 đến 5)};

    % in_first & in_last
    \node[sig_name] at (1.2, 2.3) {in\_first};
    \draw[val_line] (1.5, 1.8) -- (1.5, 2.2) -- (3.5, 2.2) -- (3.5, 1.8) -- (15.5, 1.8);
    \node[annot, text=green!50!black] at (2.5, 2.2) {Khung 0};

    \node[sig_name] at (1.2, 1.1) {in\_last};
    \draw[val_line] (1.5, 0.6) -- (11.5, 0.6) -- (11.5, 1.0) -- (13.5, 1.0) -- (13.5, 0.6) -- (15.5, 0.6);
    \node[annot, text=green!50!black] at (12.5, 1.0) {Khung 5};

    % in_f0: Vẽ bus 2 đường biên với giá trị nằm trong khoảng trắng không bị gạch ngang
    \node[sig_name] at (1.2, 0.0) {in\_f0};
    \foreach \x/\val in {1/120, 2/122, 3/244, 4/121, 5/0, 6/0} {
        \draw[data_bus, fill=white] (\x*2.0 - 0.5, -0.25) rectangle (\x*2.0 + 1.5, 0.25);
        \node[font=\scriptsize\bfseries] at (\x*2.0 + 0.5, 0.0) {\val};
    }
    \node[font=\scriptsize, text=black!40] at (14.5, 0.0) {---};

    % Chú thích lặp biên tại chu kỳ T1
    \node[annot, text=blue!80!black] at (2.5, -0.6) {Lặp biên: $z_0=z_1=z_2=120$};
    \draw[draw=blue!80!black, thick, ->] (2.5, -0.28) -- (2.5, -0.48);
    \draw[draw=blue!80!black, thick, ->] (2.5, -0.72) -- (2.5, -0.92);

    % -------------------------------------------------------------------------
    % CÁC THANH GHI NỘI TẠI
    % -------------------------------------------------------------------------
    % f0_z1 (mẫu trung tâm đang xét)
    \node[sig_name] at (1.2, -1.2) {f0\_z1};
    \foreach \x/\val in {1/120, 2/122, 3/244, 4/121, 5/0, 6/0, 7/0} {
        \draw[data_bus, fill=white] (\x*2.0 - 0.5, -1.45) rectangle (\x*2.0 + 1.5, -0.95);
        \node[font=\scriptsize\bfseries, text=blue!80!black] at (\x*2.0 + 0.5, -1.2) {\val};
    }

    % f0_z2 (mẫu quá khứ)
    \node[sig_name] at (1.2, -2.2) {f0\_z2};
    \foreach \x/\val in {1/120, 2/120, 3/122, 4/244, 5/121, 6/0, 7/0} {
        \draw[data_bus, fill=white] (\x*2.0 - 0.5, -2.45) rectangle (\x*2.0 + 1.5, -1.95);
        \node[font=\scriptsize\bfseries, text=purple!80!black] at (\x*2.0 + 0.5, -2.2) {\val};
    }

    % -------------------------------------------------------------------------
    % CÁC TÍN HIỆU ĐẦU RA
    % -------------------------------------------------------------------------
    % out_valid
    \node[sig_name] at (1.2, -3.1) {out\_valid};
    \draw[val_line] (1.5, -3.3) -- (3.5, -3.3) -- (3.5, -2.9) -- (15.5, -2.9) -- (15.5, -3.3);
    \node[annot, text=green!50!black] at (2.5, -3.1) {Chờ (0)};
    \node[annot, text=green!50!black] at (9.5, -2.9) {Mức 1 (Xuất đủ 6 khung)};

    % f0_med / out_f0
    \node[sig_name] at (1.2, -4.2) {out\_f0};
    \node[font=\scriptsize, text=black!40] at (2.5, -4.2) {---};
    \foreach \x/\val/\fr in {2/120/khung 0, 3/122/khung 1, 4/122/khung 2, 5/121/khung 3, 6/0/khung 4, 7/0/khung 5} {
        \draw[draw=orange!80!black, fill=orange!10, thick] (\x*2.0 - 0.5, -4.45) rectangle (\x*2.0 + 1.5, -3.95);
        \node[font=\scriptsize\bfseries, text=orange!90!black] at (\x*2.0 + 0.5, -4.2) {\val};
        \node[font=\scriptsize, text=black!70] at (\x*2.0 + 0.5, -4.8) {\fr};
    }

    % Mũi tên nhịp xả Flush Beat tại chu kỳ T7
    \draw[draw=red!80!black, thick, ->] (14.5, -2.5) -- node[right=3pt, font=\scriptsize\bfseries, text=red!80!black, fill=white, inner sep=1.5pt] {Flush beat!} (14.5, -3.8);

\end{tikzpicture}%
}
    \caption{Giản đồ định thời nhân quả của module \texttt{ch11\_median3.v} cho chuỗi 6 khung $[120, 122, 244, 121, 0, 0]$. Khung $k-1$ được xuất tại chu kỳ nhận khung $k$; chu kỳ T1 nạp lặp biên; chu kỳ T7 là nhịp xả Flush Beat xuất khung 5 kèm cờ \texttt{out\_last}.}
    \label{fig:ch11_median_timing}
\end{figure}

\textbf{Đọc hình.} Hình \ref{fig:ch11_median_timing} trình bày diễn biến tín hiệu theo từng chu kỳ xung nhịp từ T1 đến T7. Tại chu kỳ T1, tín hiệu $\texttt{in\_first}=1$ kích hoạt việc nạp khung đầu tiên ($120$), lặp mẫu vào cả ba thanh ghi: sau \texttt{posedge} T1 các thanh ghi $\texttt{f0\_z0}, \texttt{f0\_z1}, \texttt{f0\_z2}$ đều lưu $120$, ngõ ra $\texttt{out\_valid}=0$. Từ chu kỳ T2 đến T6, các khung mới liên tục được nạp vào, thanh ghi $\texttt{f0\_med}$ xuất ra kết quả trung vị của khung trước đó ($\texttt{out\_valid}=1$). Sau \texttt{posedge} T2, $\texttt{f0\_z2}$ giữ $120$, $\texttt{f0\_z1}$ giữ $122$, $\texttt{f0\_med}$ chốt $120$ (xuất khung 0). Sau \texttt{posedge} T3, $\texttt{f0\_z2}$ giữ $122$, $\texttt{f0\_z1}$ giữ $244$, $\texttt{f0\_med}$ chốt $122$ (xuất khung 1). Sau \texttt{posedge} T4, $\texttt{f0\_z2}$ giữ $244$, $\texttt{f0\_z1}$ giữ $121$, $\texttt{f0\_med}$ chốt $122$ (xuất khung 2, xóa gai đột biến $244$). Sau \texttt{posedge} T5, $\texttt{f0\_z2}$ giữ $121$, $\texttt{f0\_z1}$ giữ $0$, $\texttt{f0\_med}$ chốt $121$ (xuất khung 3). Sau \texttt{posedge} T6, cờ $\texttt{in\_last}=1$ báo hiệu khung 5 ($0$) đã vào, $\texttt{f0\_z2}$ và $\texttt{f0\_z1}$ đều chuyển về $0$, $\texttt{f0\_med}$ chốt $0$ (xuất khung 4). Tại chu kỳ T7 (nhịp xả Flush beat), không có mẫu vào mới, mạch tự động tính trung vị của $\{0, 0, 0\}$, xuất ra kết quả của khung cuối cùng ($0$) kèm cờ $\texttt{out\_last}=1$; các thanh ghi duy trì giá trị $0$.

\subsubsection*{Bảng vết thanh ghi từng chu kỳ (Trace Table)}
Bảng \ref{tab:ch11_trace} ghi nhận chính xác trạng thái của từng thanh ghi và tín hiệu giao tiếp trong suốt 7 chu kỳ xử lý của ví dụ mẫu, trích xuất nguyên văn từ nhật ký kiểm chứng \texttt{verify\_v3/ch11\_median\_math.log}:

\begin{table}[htbp]
\centering
\footnotesize
\setlength{\tabcolsep}{3pt}
\caption{Bảng vết thanh ghi từng chu kỳ của module \texttt{ch11\_median3.v} trên dãy 6 khung $[120, 122, 244, 121, 0, 0]$ (trạng thái thanh ghi ghi nhận sau cạnh lên xung nhịp \texttt{posedge clk}).}
\label{tab:ch11_trace}
\begin{tabular}{cccccccc>{\raggedright\arraybackslash}p{4.8cm}}
\toprule
\textbf{Chu kỳ} & \texttt{in\_valid} & \texttt{in\_f0} & \texttt{f0\_z2} & \texttt{f0\_z1} & \texttt{f0\_z0} & \texttt{out\_valid} & \texttt{f0\_med} & \textbf{Hành vi vi kiến trúc phần cứng} \\
\midrule
T1 & 1 & 120 & 120 & 120 & 120 & 0 & 0 & Nạp khung 0, lặp mẫu biên $\texttt{z0}=\texttt{z1}=\texttt{z2}=120$ \\
T2 & 1 & 122 & 120 & 122 & 122 & 1 & \textbf{120} & Xuất khung 0: $\text{med}(120, 120, 122) = 120$ \\
T3 & 1 & 244 & 122 & 244 & 244 & 1 & \textbf{122} & Xuất khung 1: $\text{med}(120, 122, 244) = 122$ \\
T4 & 1 & 121 & 244 & 121 & 121 & 1 & \textbf{122} & Xuất khung 2: $\text{med}(122, 244, 121) = 122$ (xóa gai 244!) \\
T5 & 1 & 0   & 121 & 0   & 0   & 1 & \textbf{121} & Xuất khung 3: $\text{med}(244, 121, 0) = 121$ \\
T6 & 1 & 0   & 0   & 0   & 0   & 1 & \textbf{0}   & Xuất khung 4: $\text{med}(121, 0, 0) = 0$ (\texttt{in\_last}=1) \\
T7 & 0 & --- & 0   & 0   & 0   & 1 & \textbf{0}   & Nhịp xả Flush: xuất khung 5 ($\texttt{out\_last}=1$) \\
\bottomrule
\end{tabular}
\end{table}

\textbf{Đọc bảng.} Bảng \ref{tab:ch11_trace} xác nhận mạch phần cứng chỉ mất đúng $N + 1 = 6 + 1 = 7\text{ chu kỳ xung nhịp}$ để xử lý và xả trọn vẹn $6$ khung dữ liệu. Dãy ngõ ra thu được tại các chu kỳ T2 đến T7 là $[120, 122, 122, 121, 0, 0]$, trùng khớp $100\%$ từng phần tử với mô hình toán học phần mềm.

\textit{Tóm lại: Bằng quy ước độ trễ nhân quả 1 khung (xuất khung $k-1$ tại chu kỳ $k$), kết hợp cơ chế nạp lặp biên ở chu kỳ đầu và nhịp xả tự động Flush beat ở chu kỳ cuối, phần cứng FPGA tái lập trọn vẹn kết quả của bộ lọc trung vị căn giữa.}

% ============================================================================
% 11.5. TUYẾN THIẾT KẾ VI KIẾN TRÚC HDL TỰ VIẾT
% ============================================================================
\subsection*{11.5. Tuyến thiết kế vi kiến trúc HDL \texttt{ch11\_median3.v} và lý do không dùng Intel IP}
\addcontentsline{toc}{subsection}{11.5. Tuyến thiết kế vi kiến trúc HDL \texttt{ch11\_median3.v} và lý do không dùng Intel IP}

\subsubsection*{Vấn đề thực tế: Vì sao không sử dụng thư viện Intel IP Megafunction?}
Trong các chương về lọc số (như Chương 3 với bộ lọc FIR 31 tap), ta đã sử dụng khối IP \texttt{FIR Compiler II} của Intel. Tuy nhiên, tại khối lọc trung vị Median-3 của Chương 11, câu hỏi đặt ra là: \emph{Có khối Intel Megafunction IP nào được thiết kế sẵn cho bộ lọc trung vị hay không?}

Câu trả lời là: \textbf{Không có bất kỳ khối IP chuẩn nào của Intel dành riêng cho bộ lọc trung vị}.
Các thư viện IP của Intel tập trung vào hai nhóm chức năng chính:
\begin{enumerate}
    \item Các khối số học cơ bản: bộ nhân \texttt{LPM\_MULT}, bộ chia \texttt{LPM\_DIVIDE}, bộ tích lũy \texttt{ALTERA\_MULT\_ADD}.
    \item Các bộ xử lý tín hiệu số tuyến tính: \texttt{FIR Compiler II}, \texttt{FFT Intel FPGA IP}, \texttt{NCO}.
\end{enumerate}
Bộ lọc trung vị là một giải thuật \textbf{phi tuyến tính dựa trên quan hệ thứ tự sắp xếp} (\emph{non-linear order-statistic filter}). Nếu cố tình xây dựng khối này từ các IP cơ bản của Intel, ta sẽ phải khởi tạo 3 khối so sánh độc quyền \texttt{LPM\_COMPARE}, kết nối qua nhiều tầng logic keo và tạo wrapper phức tạp. Điều này làm tăng độ cồng kềnh của dự án mà không đem lại bất kỳ ưu thế nào về diện tích hay định thời.

Ngược lại, một module Verilog tự viết cho Median-3 chỉ gói gọn trong \textbf{dưới 150 dòng mã nguồn}, hoàn toàn độc lập với các thư viện đóng kín, có thể tổng hợp tối ưu trên mọi dòng chip FPGA và mô phỏng tự do trên các công cụ mã nguồn mở như Icarus Verilog. Vì vậy, tài liệu này quyết định: \emph{Khối Median-3 được hiện thực hóa duy nhất bằng tuyến HDL tự viết}.

\subsubsection*{Sơ đồ vi kiến trúc phần cứng của module \texttt{ch11\_median3.v}}
Hình \ref{fig:ch11_median3_arch} mô tả chi tiết sơ đồ khối phần cứng của module \texttt{ch11\_median3.v}.

\begin{figure}[htbp]
    \centering
    % chuong_v3/fig_ch11_median3_arch.tex
% Sơ đồ vi kiến trúc phần cứng của module ch11_median3.v
% Bao gồm: Chuỗi thanh ghi trễ f0_z0/z1/z2, mạng 3 bộ so sánh và MUX trung vị, khối FSM điều khiển nhịp xả
\resizebox{\linewidth}{!}{%
\begin{tikzpicture}[
    >=Stealth,
    scale=0.86,
    block/.style={rectangle, draw=blue!80!black, fill=blue!5, rounded corners=2pt, minimum width=2.2cm, minimum height=1.0cm, align=center, font=\scriptsize},
    reg_box/.style={rectangle, draw=black!85, fill=gray!25, rounded corners=2pt, minimum width=1.4cm, minimum height=1.1cm, align=center, font=\scriptsize\bfseries},
    fsm_box/.style={rectangle, draw=purple!80!black, fill=purple!5, rounded corners=2pt, minimum width=2.8cm, minimum height=1.4cm, align=center, font=\scriptsize},
    wire/.style={draw=black!80, thick, ->},
    ctrl_wire/.style={draw=purple!80!black, thick, dashed, ->},
    bus_lbl/.style={font=\scriptsize, fill=white, inner sep=1.2pt}
]

    % -------------------------------------------------------------------------
    % KHỐI ĐIỀU KHIỂN FSM
    % -------------------------------------------------------------------------
    \node[fsm_box] (fsm) at (2.0, 4.2) {Khối FSM điều khiển\\(IDLE / STREAM / FLUSH)};
    \draw[wire] (-1.2, 4.6) -- node[above=2pt, bus_lbl] {\texttt{in\_valid}, \texttt{in\_first}} (fsm.170);
    \draw[wire] (-1.2, 3.8) -- node[below=2pt, bus_lbl] {\texttt{in\_last}} (fsm.190);

    % -------------------------------------------------------------------------
    % ĐƯỜNG DỮ LIỆU: CHUỖI THANH GHI TRỄ
    % -------------------------------------------------------------------------
    \node[reg_box] (rz0) at (2.0, 1.8) {Thanh ghi\\$\texttt{f0\_z0}$\\(16-bit)};
    \node[reg_box] (rz1) at (5.2, 1.8) {Thanh ghi\\$\texttt{f0\_z1}$\\(16-bit)};
    \node[reg_box] (rz2) at (8.2, 1.8) {Thanh ghi\\$\texttt{f0\_z2}$\\(16-bit)};

    % Đường nạp ngõ vào in_f0
    \draw[wire] (-1.2, 1.8) -- node[above=2pt, bus_lbl] {\texttt{in\_f0} [16]} (rz0.west);
    \draw[wire] (rz0.east) -- node[above=4pt, font=\scriptsize\bfseries, fill=white, inner sep=1.5pt] {shift} (rz1.west);
    \draw[wire] (rz1.east) -- node[above=4pt, font=\scriptsize\bfseries, fill=white, inner sep=1.5pt] {shift} (rz2.west);

    % Tín hiệu điều khiển lặp mẫu biên và nhịp xả từ FSM
    \draw[ctrl_wire] (fsm.south) -- node[right=4pt, font=\scriptsize, fill=white, inner sep=1.5pt, text=purple!80!black] {en, lặp biên} (rz0.north);

    % -------------------------------------------------------------------------
    % MẠNG 3 BỘ SO SÁNH VÀ GHÉP KÊNH TRUNG VỊ
    % -------------------------------------------------------------------------
    \node[block] (cmp_net) at (11.2, 1.8) {Mạng 3 bộ so sánh:\\$a \le b,\ b \le c,\ a \le c$\\($a=z_2, b=z_1, c=z_0$)};
    \node[block] (mux_med) at (16.0, 1.8) {Bộ ghép kênh MUX\\chọn Trung vị 16-bit};

    % Đường nối từ thanh ghi vào mạng so sánh (chạy bên dưới thanh ghi, khoảng cách >= 3 mm)
    \draw[wire] (rz2.east) -- (cmp_net.west);
    \draw[wire] (rz1.south) -- ++(0, -0.65) -- (9.9, 0.60) |- (cmp_net.205);
    \draw[wire] (rz0.south) -- ++(0, -1.05) -- (9.6, 0.20) |- (cmp_net.220);

    \draw[wire] (cmp_net.east) -- node[above=5pt, font=\scriptsize, fill=white, inner sep=2pt, rounded corners=1pt] {3 cờ so sánh} (mux_med.west);

    % -------------------------------------------------------------------------
    % THANH GHI CHỐT NGÕ RA f0_med
    % -------------------------------------------------------------------------
    \node[reg_box] (rmed) at (19.6, 1.8) {Thanh ghi\\$\texttt{f0\_med}$\\(16-bit)};
    \draw[wire] (mux_med.east) -- (rmed.west);
    \draw[wire] (rmed.east) -- node[above=2pt, bus_lbl] {\texttt{out\_f0} [16]} (22.2, 1.8);

    % Tín hiệu out_valid và out_last từ FSM
    \draw[ctrl_wire] (fsm.east) -- ++(14.2, 0) -- node[above=2pt, bus_lbl] {\texttt{out\_valid}, \texttt{out\_last}} (22.2, 4.2);

    % -------------------------------------------------------------------------
    % KHUNG BAO VÀ CHÚ THÍCH
    % -------------------------------------------------------------------------
    \draw[draw=black!40, dashed, rounded corners=3pt] (0.8, -0.15) rectangle (20.8, 2.7);
    \node[above=2pt, font=\scriptsize\bfseries, text=black!60] at (11.0, 2.7) {Đường truyền dữ liệu Median-3 (1 chu kỳ xung nhịp)};

\end{tikzpicture}%
}

    \caption{Sơ đồ vi kiến trúc phần cứng của module \texttt{ch11\_median3.v}. Chuỗi thanh ghi dịch $\texttt{f0\_z0} \to \texttt{f0\_z1} \to \texttt{f0\_z2}$ duy trì 3 mẫu quan sát; mạng 3 bộ so sánh nhị phân điều khiển MUX chọn trung vị; khối FSM điều phối trạng thái dòng chảy và nhịp xả tự động.}
    \label{fig:ch11_median3_arch}
\end{figure}

\textbf{Đọc hình.} Hình \ref{fig:ch11_median3_arch} thể hiện rõ sự phân tách giữa đường dữ liệu (\emph{datapath}) và khối điều khiển (\emph{control path}):
\begin{itemize}
    \item \textbf{Đường dữ liệu:} Mẫu ngõ vào $\texttt{in\_f0}$ (16-bit) được đưa vào chuỗi thanh ghi trễ gồm $\texttt{f0\_z0}$, $\texttt{f0\_z1}$ và $\texttt{f0\_z2}$. Mạng 3 bộ so sánh tổ hợp ($a \le b$, $b \le c$, $a \le c$) tiếp nhận dữ liệu từ các thanh ghi và xuất ra 3 cờ điều khiển cho bộ ghép kênh MUX. Giá trị trung vị sau MUX được chốt trực tiếp vào thanh ghi ngõ ra $\texttt{f0\_med}$.
    \item \textbf{Khối điều khiển FSM:} Nhận các tín hiệu bắt tay $\texttt{in\_valid}$, $\texttt{in\_first}$ và $\texttt{in\_last}$. FSM điều khiển việc nạp lặp biên tại trạng thái \texttt{S\_IDLE}, duy trì dòng dịch dữ liệu tại \texttt{S\_STREAM}, và tự động chuyển sang \texttt{S\_FLUSH} để phát nhịp xả khi phát hiện khung cuối cùng.
\end{itemize}

\subsubsection*{Mã nguồn Verilog hiện thực hóa module \texttt{ch11\_median3.v}}
Listing \ref{lst:ch11_median3_v} trình bày toàn bộ phần thân tính toán của module \texttt{ch11\_median3.v} (từ dòng 40 đến dòng 160):

\lstinputlisting[language=Verilog, firstline=40, lastline=160, firstnumber=40, caption={Mã nguồn Verilog hiện thực hóa vi kiến trúc bộ lọc trung vị \texttt{ch11\_median3.v}.}, label={lst:ch11_median3_v}]{rtl/v3/ch11_median3.v}

\subsubsection*{Phân tích chi tiết từng khối mã nguồn Verilog}
\begin{itemize}
    \item \textbf{Dòng 40--54 (Khai báo FSM và Thanh ghi):} Định nghĩa 3 trạng thái điều khiển: \texttt{S\_IDLE} ($0$), \texttt{S\_STREAM} ($1$), và \texttt{S\_FLUSH} ($2$). Bốn thanh ghi 16-bit chuẩn tắc được khai báo: $\texttt{f0\_z0}$ (mẫu mới), $\texttt{f0\_z1}$ (mẫu trung tâm), $\texttt{f0\_z2}$ (mẫu quá khứ) và $\texttt{f0\_med}$ (thanh ghi chốt trung vị ngõ ra). Cổng ngõ ra $\texttt{out\_f0}$ được nối trực tiếp với $\texttt{f0\_med}$ bằng lệnh gán liên tục \texttt{assign out\_f0 = f0\_med}.
    \item \textbf{Dòng 60--66 (Mạng so sánh tổ hợp):} Ba toán hạng so sánh được gán: $\texttt{med\_a} = \texttt{f0\_z2}$, $\texttt{med\_b} = \texttt{f0\_z1}$, và $\texttt{med\_c} = (\texttt{state} == \texttt{S\_FLUSH})\ ?\ \texttt{f0\_z1} : \texttt{in\_f0}$. Điểm đặc sắc ở đây là \textbf{chia sẻ tài nguyên bộ so sánh}: trong nhịp xả \texttt{S\_FLUSH}, toán hạng thứ ba tự động chuyển sang đọc mẫu cuối $\texttt{f0\_z1}$ mà không cần tạo thêm bất kỳ bộ so sánh phụ nào! Ba dây cờ \texttt{cmp\_ab}, \texttt{cmp\_bc}, \texttt{cmp\_ac} được tính toán đồng thời bằng phép so sánh không dấu 16-bit (\texttt{<=}).
    \item \textbf{Dòng 68--85 (Bộ ghép kênh MUX trung vị):} Khối \texttt{always @(*)} hiện thực hóa chính xác cây quyết định \eqref{eq:median3_cases}. Kết quả trung vị được đưa vào biến tổ hợp \texttt{med\_comb}.
    \item \textbf{Dòng 90--99 (Reset bất đồng bộ):} Khi có xung $\texttt{!rst\_n}$, toàn bộ FSM và thanh ghi bị xóa về $0$, các cờ ngõ ra $\texttt{out\_valid}$ và $\texttt{out\_last}$ bị kéo xuống mức $0$.
    \item \textbf{Dòng 100--116 (Trạng thái S\_IDLE):} Chờ tín hiệu $\texttt{in\_valid \&\& in\_first}$. Khi nhận được khung 0, giá trị $\texttt{in\_f0}$ được nạp đồng thời vào cả $\texttt{z0}, \texttt{z1}, \texttt{z2}$ để tạo biên lặp đối xứng. Nếu tệp đặc biệt chỉ có đúng 1 khung duy nhất ($\texttt{in\_last} = 1$), FSM nhảy thẳng sang \texttt{S\_FLUSH}; ngược lại chuyển sang \texttt{S\_STREAM}.
    \item \textbf{Dòng 118--141 (Trạng thái S\_STREAM):} Mỗi khi có mẫu hợp lệ $\texttt{in\_valid} = 1$: chốt $\texttt{f0\_med} \le \texttt{med\_comb}$, bật $\texttt{out\_valid} = 1$, đồng thời dịch chuyển thanh ghi $\texttt{f0\_z2} \le \texttt{f0\_z1}$ và $\texttt{f0\_z1} \le \texttt{in\_f0}$. Nếu gặp $\texttt{in\_last} = 1$, FSM chuyển tiếp sang \texttt{S\_FLUSH}. Nếu không có mẫu mới, giữ nguyên giá trị thanh ghi và hạ $\texttt{out\_valid} = 0$.
    \item \textbf{Dòng 143--151 (Nhịp xả \texttt{S\_FLUSH}):} Chốt kết quả mẫu lặp cuối vào $\texttt{f0\_med}$, bật đồng thời $\texttt{out\_valid} = 1$ và $\texttt{out\_last} = 1$, rồi trở về \texttt{S\_IDLE}.
\end{itemize}

\textit{Tóm lại: Module \texttt{ch11\_median3.v} được thiết kế tinh gọn dưới 150 dòng Verilog thuần túy; kỹ thuật chia sẻ mạng so sánh giữa dòng chảy chính và nhịp xả giúp mạch tiêu hao 0 khối DSP và 0 khối M10K, đáp ứng thông lượng 1 khung/chu kỳ.}

% ============================================================================
% 11.6. MÔ HÌNH SONG SINH PYTHON TWIN VÀ KIỂM CHỨNG
% ============================================================================
\subsection*{11.6. Mô hình song sinh Python Twin và kịch bản kiểm chứng mô phỏng}
\addcontentsline{toc}{subsection}{11.6. Mô hình song sinh Python Twin và kịch bản kiểm chứng mô phỏng}

\subsubsection*{Cấu trúc mô hình song sinh phần mềm (Python Twin)}
Để bảo đảm mã RTL Verilog hoạt động chính xác theo từng chu kỳ xung nhịp và từng bit dữ liệu, mô hình song sinh \texttt{Median3Twin} được xây dựng trong tệp \texttt{verify\_v3/ch11\_twin.py}.
Hàm cập nhật chu kỳ xung nhịp \texttt{clock()} của mô hình song sinh được trình bày trong Listing \ref{lst:ch11_twin_py} (dòng 66 đến dòng 141):

\lstinputlisting[language=Python, firstline=66, lastline=141, firstnumber=66, caption={Hàm cập nhật chu kỳ xung nhịp \texttt{clock()} của mô hình song sinh Python Twin (\texttt{verify\_v3/ch11\_twin.py}).}, label={lst:ch11_twin_py}]{verify_v3/ch11_twin.py}

\subsubsection*{Bảng đối chiếu dòng mã nguồn giữa Python Twin và Verilog HDL}
Bảng \ref{tab:ch11_twin_vs_hdl} đối chiếu chi tiết từng vị trí dòng lệnh giữa mô hình Python Twin (\texttt{verify\_v3/ch11\_twin.py}) và mã phần cứng Verilog (\texttt{rtl/v3/ch11\_median3.v}).

\begin{table}[htbp]
\centering
\small
\caption{Bảng đối chiếu từng dòng giữa Python Twin (\texttt{verify\_v3/ch11\_twin.py}) và Verilog HDL (\texttt{rtl/v3/ch11\_median3.v}).}
\label{tab:ch11_twin_vs_hdl}
\begin{tabular}{p{3.8cm}p{3.8cm}p{7.2cm}}
\toprule
\textbf{Dòng Python Twin} & \textbf{Dòng Verilog HDL} & \textbf{Ý nghĩa vi kiến trúc mạch phần cứng} \\
\midrule
Dòng 57--64 & Dòng 91--98 & Khởi tạo và Reset bất đồng bộ (\texttt{!rst\_n}), xóa sạch các thanh ghi và cờ báo về $0$. \\
Dòng 33--51 & Dòng 60--85 & Mạng logic 3 bộ so sánh nhị phân (\texttt{<=}) và cây ghép kênh MUX chọn phần tử trung vị. \\
Dòng 82--98 & Dòng 100--116 & \textbf{Trạng thái S\_IDLE:} Nạp khung đầu tiên $x[0]$, lặp mẫu biên vào $\texttt{f0\_z0}, \texttt{f0\_z1}$ và $\texttt{f0\_z2}$. \\
Dòng 100--121 & Dòng 118--141 & \textbf{Trạng thái S\_STREAM:} Xuất kết quả trung vị khung $k-1$, cập nhật thanh ghi: $\texttt{f0\_z2} \leftarrow \texttt{f0\_z1}$, $\texttt{f0\_z1} \leftarrow \texttt{in\_f0}$, $\texttt{f0\_z0} \leftarrow \texttt{in\_f0}$. \\
Dòng 122--131 & Dòng 143--151 & \textbf{Trạng thái S\_FLUSH:} Nhịp xả tự động, lặp mẫu cuối $\texttt{f0\_z1}$, xuất kết quả khung $N-1$ kèm cờ $\texttt{out\_last}=1$. \\
Dòng 133--140 & Dòng 126--128, 149 & Cập nhật trạng thái các thanh ghi tại cạnh lên xung nhịp (\texttt{posedge clk}). \\
\bottomrule
\end{tabular}
\end{table}

\textbf{Đọc bảng.} Bảng \ref{tab:ch11_twin_vs_hdl} khẳng định tính tương đồng tuyệt đối giữa mô hình phần mềm và mã phần cứng. Từng nhánh rẽ điều kiện FSM, thứ tự dịch chuyển thanh ghi và logic mạng so sánh trung vị đều phản ánh chính xác cấu trúc vi kiến trúc trên chip silicon.

\subsubsection*{Quy ước lấy mẫu Testbench hai pha (Sampling Convention)}
Để loại trừ hiện tượng tranh chấp dữ liệu (race conditions) trong mô phỏng số, Testbench trong tệp \texttt{tb/v3/tb\_ch11\_median3.v} tuân thủ nghiêm ngặt \textbf{quy ước lấy mẫu hai pha (2-phase sampling convention)}:
\begin{itemize}
    \item \textbf{Pha kích thích ngõ vào (Drive phase):} Mọi tín hiệu ngõ vào ($\texttt{in\_valid}$, $\texttt{in\_first}$, $\texttt{in\_last}$, $\texttt{in\_f0}$) được kích hoạt tại cạnh xuống của xung nhịp (\texttt{@(negedge clk)}). Khi đó tín hiệu có trọn vẹn nửa chu kỳ ($10\text{ ns}$) để ổn định trước cạnh lên.
    \item \textbf{Pha kiểm tra ngõ ra (Sample phase):} Sau khi phát hiện cạnh lên (\texttt{@(posedge clk)}), Testbench chờ thêm trễ giữ $\texttt{\#1}$ ($1\text{ ns}$), rồi mới tiến hành đọc và so sánh các ngõ ra đã chốt ($\texttt{out\_valid}$, $\texttt{out\_last}$, $\texttt{out\_f0}$) với tệp dữ liệu kỳ vọng sinh ra từ Python Twin.
\end{itemize}
Testbench không bao giờ ngắt sớm khi gặp lỗi mà duyệt hết toàn bộ tệp kích thích để thống kê đầy đủ số lỗi so sánh.

\subsubsection*{Bộ vectơ kiểm thử góc và xác minh mô phỏng}
Tệp kích thích \texttt{verify\_v3/ch11\_twin.py} đã sinh ra tổng cộng \textbf{1.065 chu kỳ mô phỏng}, kiểm tra trọn vẹn \textbf{609 kết quả ngõ ra hoàn chỉnh}. Bộ vectơ bao phủ toàn diện các kịch bản thử nghiệm:
\begin{enumerate}
    \item Ví dụ tính tay đề cương $[120, 122, 244, 121, 0, 0]$ (chạy liền mạch và chạy có chu kỳ rỗi stall ngẫu nhiên).
    \item Các trường hợp biên đặc thù: tệp chỉ có 1 khung duy nhất ($N = 1$) và tệp có 2 khung ($N = 2$).
    \item Các ca gai đơn lẻ 1 khung (nhân đôi $1600 \to 3200 \to 1600$ và chia đôi quãng tám $1600 \to 800 \to 1600$).
    \item Các ca gai kép kéo dài 2 khung liên tiếp ($1800 \to 3600 \to 3600 \to 1800$).
    \item Các chuỗi số 0 (đoạn vô thanh dài ngắn khác nhau), đốm sáng giả trong im lặng và mất tiếng đơn khung.
    \item Toàn bộ 277 khung dữ liệu thực tế từ tệp sạch \texttt{synth\_glide\_vuv\_16k.wav}.
    \item Toàn bộ 277 khung dữ liệu thực tế từ tệp có nhiễu \texttt{synth\_glide\_vuv\_16k\_snr5.wav}.
\end{enumerate}

Kịch bản thực thi tự động \texttt{verify\_v3/run\_ch11\_iverilog.sh} đã được thiết lập với định dạng dòng Unix LF và cờ thực thi, sẵn sàng để người phản biện chạy trên môi trường Linux:
\begin{quote}
\footnotesize
\texttt{iverilog -g2012 -Wall -I rtl/v3 -o verify\_v3/sim\_ch11\_median3 \textbackslash}\\
\texttt{\ \ \ \ tb/v3/tb\_ch11\_median3.v rtl/v3/ch11\_median3.v}\\
\texttt{vvp verify\_v3/sim\_ch11\_median3}
\end{quote}

\todo{box: Kết quả chạy iverilog tb\_ch11\_median3 trên Linux do người phản biện cung cấp và điền vào đây.}

\subsubsection*{Kiểm chứng mã tham chiếu C trên bộ xử lý HPS}
Để phục vụ việc kiểm thử phần mềm trên bộ xử lý ARM Cortex-A9 của HPS, tệp mã nguồn C \texttt{sw\_v3/ch11\_median3.c} và chương trình kiểm thử \texttt{sw\_v3/test\_ch11\_median3.c} đã được biên dịch bằng GCC:
\begin{quote}
\footnotesize
\texttt{gcc -Wall -Wextra -O2 sw\_v3/test\_ch11\_median3.c sw\_v3/ch11\_median3.c \textbackslash}\\
\texttt{\ \ \ \ -Isw\_v3 -o test\_ch11\_median3}\\
\texttt{./test\_ch11\_median3}
\end{quote}
Kết quả thực thi xác nhận \textbf{4/4 bài kiểm tra đơn vị vượt qua thành công} (0 cảnh báo, 0 lỗi), bảo đảm tính tương đương toán học giữa mã nguồn C trên CPU và phần cứng Verilog trên FPGA.

\subsubsection*{Bảng so sánh tổng hợp giữa Tuyến Intel IP và Tuyến HDL tự viết}
Bảng \ref{tab:ch11_ip_vs_hdl} tổng hợp so sánh giữa giải pháp dùng Intel IP (giả định) và module HDL tự viết \texttt{ch11\_median3.v}.

\begin{table}[htbp]
\centering
\small
\caption{Bảng so sánh tổng hợp giữa Tuyến Intel IP (giả định) và Tuyến HDL tự viết (\texttt{ch11\_median3.v}).}
\label{tab:ch11_ip_vs_hdl}
\begin{tabular}{p{3.5cm}p{5.5cm}p{5.8cm}}
\toprule
\textbf{Tiêu chí kỹ thuật} & \textbf{Tuyến (A): Thư viện Intel IP (Giả định)} & \textbf{Tuyến (B): HDL tự viết (\texttt{ch11\_median3.v})} \\
\midrule
Khả năng hỗ trợ sẵn có & Không có IP chuyên dụng cho lọc trung vị; phải ghép 3 khối \texttt{LPM\_COMPARE}. & Thiết kế Verilog-2001 hoàn chỉnh, chuyên biệt, dưới 150 dòng mã. \\
Tài nguyên M10K RAM & 0 khối M10K. & \textbf{0 khối M10K} (hoàn toàn dùng thanh ghi ALM). \\
Tài nguyên DSP Blocks & 0 khối DSP. & \textbf{0 khối DSP} (phép so sánh không dấu). \\
Tài nguyên Logic ALM & Tốn thêm logic keo wrapper IP (phải đo trong Quartus). & Chỉ gồm 4 thanh ghi 16-bit và 3 bộ so sánh (phải đo trong Quartus). \\
Tần số cực đại $F_{\max}$ & Phụ thuộc định tuyến giữa các wrapper (phải đo trong Quartus). & Mạch tổ hợp cực ngắn, đáp ứng xung nhịp $50\text{ MHz}$ (phải đo trong Quartus). \\
Độ trễ xử lý (Latency) & Phụ thuộc cấu hình wrapper IP. & \textbf{Đúng 1 chu kỳ xung nhịp} (sau khi có đủ mẫu). \\
Độ trễ nhân quả (Causal) & 1 khung thời gian ($10\text{ ms}$). & \textbf{Đúng 1 khung thời gian} (xuất khung $k-1$ tại chu kỳ $k$). \\
Xử lý nhịp xả (Flush) & Phải tự thiết kế thêm logic ngoại vi. & Tự động sinh nhịp xả bằng FSM nội tại kèm cờ \texttt{out\_last}. \\
Tính khả chuyển mô phỏng & Đòi hỏi thư viện \texttt{220model} của Quartus. & Độc lập hoàn toàn, mô phỏng tự do trên Icarus Verilog, Verilator, ModelSim. \\
\bottomrule
\end{tabular}
\end{table}

\textit{Tóm lại: Module \texttt{ch11\_median3.v} và mô hình Python Twin đạt độ tương thích chu kỳ-đối-chu kỳ trên 1.065 chu kỳ mô phỏng; thiết kế HDL tự viết tối ưu hóa triệt để diện tích silicon (0 DSP, 0 M10K) và giải quyết trọn vẹn bài toán nhịp xả ở khung cuối.}

% ============================================================================
% 11.7. BÀI TẬP CUỐI CHƯƠNG CÓ LỜI GIẢI CHI TIẾT
% ============================================================================
\subsection*{11.7. Bài tập cuối chương có lời giải chi tiết}
\addcontentsline{toc}{subsection}{11.7. Bài tập cuối chương có lời giải chi tiết}

\paragraph{Bài tập 11.1 (Dẫn xuất toán học sự suy biến của Viterbi thành Per-Frame Argmax).}
Giả sử một tệp âm thanh có $N$ khung, mỗi khung có $M$ ứng viên cao độ với điểm số cục bộ $\text{local\_cost}[t][i]$ ($t = 1..N, i = 1..M$).
\begin{enumerate}[label=(\alph*)]
    \item Viết lại phương trình quy hoạch động Viterbi của Praat khi chi phí chuyển trạng thái $\text{transitionCost} = 0$.
    \item Chứng minh bằng quy nạp toán học rằng đường đi tối ưu Viterbi toàn cục trùng khớp hoàn toàn với phép chọn cực đại độc lập từng khung: $\text{place}_t = \arg\max_{i} \text{local\_cost}[t][i]$.
    \item Nêu ý nghĩa của định lý này trong việc thiết lập cấu hình đối chuẩn Praat R2.
\end{enumerate}

\textit{Lời giải chi tiết:}
\begin{enumerate}[label=(\alph*)]
    \item Khi $\text{transitionCost}(i_1 \to i_2) = 0$, điểm số tích lũy tại khung $t$ là:
    \begin{equation*}
        \delta[t][i_2] = \max_{i_1 \in [1, M]} \left( \delta[t-1][i_1] - 0 + \text{local\_cost}[t][i_2] \right)
    \end{equation*}
    Vì số hạng $\text{local\_cost}[t][i_2]$ không phụ thuộc vào biến chỉ số $i_1$, ta tách nó ra khỏi toán tử cực đại:
    \begin{equation*}
        \delta[t][i_2] = \left( \max_{i_1 \in [1, M]} \delta[t-1][i_1] \right) + \text{local\_cost}[t][i_2]
    \end{equation*}
    Con trỏ truy vết ngược là:
    \begin{equation*}
        \psi[t][i_2] = \arg\max_{i_1 \in [1, M]} \delta[t-1][i_1]
    \end{equation*}
    \item Đặt $M_{t-1} = \max_{i_1} \delta[t-1][i_1]$ là một hằng số đối với mọi $i_2$.
    Tại khung cuối cùng $N$:
    \begin{equation*}
        \text{place}_N = \arg\max_{i_2} \delta[N][i_2] = \arg\max_{i_2} \left( M_{N-1} + \text{local\_cost}[N][i_2] \right) = \arg\max_{i_2} \text{local\_cost}[N][i_2]
    \end{equation*}
    Khi truy vết ngược về khung $N-1$:
    \begin{equation*}
        \text{place}_{N-1} = \psi[N][\text{place}_N] = \arg\max_{i_1} \delta[N-1][i_1] = \arg\max_{i_1} \text{local\_cost}[N-1][i_1]
    \end{equation*}
    Lặp lại tương tự cho mọi khung $t = N-2, \dots, 1$, ta luôn có:
    \begin{equation*}
        \text{place}_t = \arg\max_{i} \text{local\_cost}[t][i]
    \end{equation*}
    Đường đi toàn cục tối ưu trùng khớp hoàn toàn với phép chọn cực đại độc lập từng khung.
    \item Ý nghĩa: Bằng cách tắt chi phí chuyển trạng thái ($\text{jump} = 0, \text{vuv} = 0$), cấu hình Praat R2 đưa bài toán Viterbi toàn cục về đúng bản chất tính toán đơn khung của phần cứng FPGA trước median, tạo ra một chuẩn đối sánh ngang hàng trực tiếp (\emph{apples-to-apples}) đạt độ tương thích toán học tuyệt đối (nhật ký \texttt{verify\_v3/ch11\_exercises.log}).
\end{enumerate}

\paragraph{Bài tập 11.2 (Tính toán thủ công thu nhỏ trên dãy 10 khung thử nghiệm).}
Cho một dãy 10 khung cao độ ngõ vào:
\begin{equation*}
    \mathbf{x = [0,\, 150,\, 300,\, 150,\, 155,\, 310,\, 310,\, 155,\, 0,\, 160]}
\end{equation*}
\begin{enumerate}[label=(\alph*)]
    \item Hãy viết dãy sau khi đã đệm lặp mẫu ở hai biên theo quy tắc chuẩn.
    \item Tính giá trị ngõ ra của bộ lọc trung vị Median-3 tại từng khung từ $k = 0$ đến $k = 9$.
    \item Chỉ ra khung nào là gai đơn lẻ được sửa thành công, khung nào là gai kép không sửa được, và khung nào là hiện tượng mất tiếng được hàn gắn.
\end{enumerate}

\textit{Lời giải chi tiết:}
\begin{enumerate}[label=(\alph*)]
    \item Dãy sau khi đệm lặp mẫu biên ($x_0 = 0$ ở đầu và $x_9 = 160$ ở cuối):
    \begin{equation*}
        x_{\text{pad}} = [0,\, \mathbf{0,\, 150,\, 300,\, 150,\, 155,\, 310,\, 310,\, 155,\, 0,\, 160},\, 160]
    \end{equation*}
    \item Tính toán trung vị 3 điểm trên từng cửa sổ:
    \begin{itemize}
        \item Khung $0$: cửa sổ $[0, 0, 150] \implies \text{med} = \mathbf{0}$.
        \item Khung $1$: cửa sổ $[0, 150, 300] \implies \text{med} = \mathbf{150}$.
        \item Khung $2$: cửa sổ $[150, 300, 150] \implies \text{med} = \mathbf{150}$.
        \item Khung $3$: cửa sổ $[300, 150, 155] \implies \text{med} = \mathbf{155}$.
        \item Khung $4$: cửa sổ $[150, 155, 310] \implies \text{med} = \mathbf{155}$.
        \item Khung $5$: cửa sổ $[155, 310, 310] \implies \text{med} = \mathbf{310}$.
        \item Khung $6$: cửa sổ $[310, 310, 155] \implies \text{med} = \mathbf{310}$.
        \item Khung $7$: cửa sổ $[310, 155, 0] \implies \text{med} = \mathbf{155}$.
        \item Khung $8$: cửa sổ $[155, 0, 160] \implies \text{med} = \mathbf{155}$.
        \item Khung $9$: cửa sổ $[0, 160, 160] \implies \text{med} = \mathbf{160}$.
    \end{itemize}
    Dãy kết quả ngõ ra là:
    \begin{equation*}
        \mathbf{y = [0,\, 150,\, 150,\, 155,\, 155,\, 310,\, 310,\, 155,\, 155,\, 160]}
    \end{equation*}
    \item Phân tích các trường hợp:
    \begin{itemize}
        \item \textbf{Khung 2:} Gai đơn $300\text{ Hz}$ giữa hai khung $150\text{ Hz}$ đã được triệt tiêu thành công về $150\text{ Hz}$.
        \item \textbf{Khung 5 và 6:} Gai kép $310\text{ Hz}$ kéo dài liên tiếp 2 khung không thể sửa được, ngõ ra vẫn giữ nguyên giá trị lỗi $310\text{ Hz}$ do giá trị này chiếm $2/3$ phần tử trong cửa sổ.
        \item \textbf{Khung 8:} Khung mất tiếng đột ngột $0\text{ Hz}$ giữa hai khung hữu thanh $155\text{ Hz}$ và $160\text{ Hz}$ đã được hàn gắn thành công thành $155\text{ Hz}$ (nhật ký \texttt{verify\_v3/ch11\_exercises.log}).
    \end{itemize}
\end{enumerate}

\paragraph{Bài tập 11.3 (Phân tích định thời nhân quả phần cứng và nhịp xả Flush Beat).}
Cho một tệp âm thanh ngắn gồm đúng 4 khung ngõ vào: $x = [100, 200, 300, 400]$.
\begin{enumerate}[label=(\alph*)]
    \item Lập bảng phân tích diễn biến của các tín hiệu \texttt{in\_valid}, \texttt{in\_first}, \texttt{in\_last}, các thanh ghi $\texttt{f0\_z0}, \texttt{f0\_z1}, \texttt{f0\_z2}$, và ngõ ra $\texttt{out\_valid}, \texttt{out\_last}, \texttt{out\_f0}$ theo từng chu kỳ xung nhịp.
    \item Chứng minh rằng cần chính xác $N + 1 = 5\text{ chu kỳ}$ để phần cứng xuất đủ 4 kết quả trùng khớp với công thức đệm biên đối xứng.
\end{enumerate}

\textit{Lời giải chi tiết:}
\begin{enumerate}[label=(\alph*)]
    \item Bảng diễn biến tín hiệu theo từng chu kỳ:
    \begin{table}[H]
    \centering
    \footnotesize
    \setlength{\tabcolsep}{2.5pt}
    \begin{tabular}{cccccccccp{4.8cm}}
    \toprule
    \textbf{Chu kỳ} & \texttt{in\_vld} & \texttt{first} & \texttt{last} & \texttt{in\_f0} & \texttt{z2} & \texttt{z1} & \texttt{out\_vld} & \texttt{out\_f0} & \textbf{Hành động vi kiến trúc} \\
    \midrule
    T1 & 1 & 1 & 0 & 100 & 100 & 100 & 0 & 0 & Nạp khung 0, lặp biên $\texttt{z1}=\texttt{z2}=100$ \\
    T2 & 1 & 0 & 0 & 200 & 100 & 200 & 1 & \textbf{100} & Xuất khung 0: $\text{med}(100, 100, 200) = 100$ \\
    T3 & 1 & 0 & 0 & 300 & 200 & 300 & 1 & \textbf{200} & Xuất khung 1: $\text{med}(100, 200, 300) = 200$ \\
    T4 & 1 & 0 & 1 & 400 & 300 & 400 & 1 & \textbf{300} & Xuất khung 2: $\text{med}(200, 300, 400) = 300$ \\
    T5 & 0 & 0 & 0 & --- & 300 & 400 & 1 & \textbf{400} & Nhịp xả: $\text{med}(300, 400, 400) = 400$ (\texttt{last}=1) \\
    \bottomrule
    \end{tabular}
    \end{table}
    \item Phân tích: Tại chu kỳ T1, khung 0 được nạp nhưng chưa thể tính do chưa có khung 1. Tại các chu kỳ T2, T3, T4, mạch lần lượt xuất các khung 0, 1, 2. Tại chu kỳ T5 (nhịp xả), mạch kích hoạt $\texttt{out\_last}=1$ và xuất nốt khung cuối cùng 3. Tổng số chu kỳ để hoàn thành $N=4$ khung là đúng $N+1 = 5\text{ chu kỳ xung nhịp}$ (nhật ký \texttt{verify\_v3/ch11\_exercises.log}).
\end{enumerate}

\paragraph{Bài tập 11.4 (Phân tích độ phức tạp mạng so sánh và tối ưu tài nguyên FPGA).}
\begin{enumerate}[label=(\alph*)]
    \item So sánh chi phí phần cứng giữa mạng so sánh 3 phần tử dùng trong module \texttt{ch11\_median3.v} với một bộ sắp xếp nổi bọt (Bubble Sort) 3 phần tử thông thường.
    \item Đánh giá tài nguyên tiêu hao (ALM, DSP, M10K) của khối Median-3 trên kiến trúc Cyclone V so với giải pháp Viterbi Path Finder.
\end{enumerate}

\textit{Lời giải chi tiết:}
\begin{enumerate}[label=(\alph*)]
    \item Mạng so sánh 3 phần tử trong module \texttt{ch11\_median3.v} thực hiện song song 3 phép so sánh không dấu 16-bit ($a \le b$, $b \le c$, $a \le c$). Mỗi phép so sánh là một chuỗi carry-chain 16-bit trong khối ALM, tạo ra 3 cờ nhị phân điều khiển MUX 3 nhánh. Thời gian lan truyền chỉ mất 1 tầng logic tổ hợp. Ngược lại, thuật toán Bubble Sort đòi hỏi nhiều chu kỳ hoán đổi (tối đa 3 bước hoán đổi) hoặc chuỗi 3 khối so sánh-hoán vị nối tiếp (Compare-and-Swap units), làm tăng độ trễ đường truyền và tốn nhiều thanh ghi trung gian hơn.
    \item Đánh giá tài nguyên trên Cyclone V:
    \begin{itemize}
        \item Khối Median-3: Tiêu tốn \textbf{0 khối DSP}, \textbf{0 khối RAM M10K}, chỉ gồm 4 thanh ghi 16-bit ($\texttt{f0\_z0}, \texttt{f0\_z1}, \texttt{f0\_z2}, \texttt{f0\_med}$) cùng mạng logic so sánh tổ hợp (số lượng ALM cụ thể phải đo trong Quartus Prime). Thông lượng đạt tuyệt đối 1 khung / chu kỳ.
        \item Giải pháp Viterbi: Cần ít nhất $60\text{ KB}$ đến vài trăm KB SRAM nội chip để lưu ma trận Trellis, tiêu tốn nhiều khối RAM M10K, đồng thời cần logic tính khoảng cách logarit phức tạp và bộ nhớ truy vết ngược phi nhân quả.
    \end{itemize}
    Do đó, Median-3 là sự lựa chọn tối ưu tuyệt đối về mặt tài nguyên và đáp ứng xuất sắc yêu cầu thời gian thực trên chip FPGA (nhật ký \texttt{verify\_v3/ch11\_exercises.log}).
\end{enumerate}

\textit{Tóm lại: Toàn bộ 4 bài tập cuối chương đã được giải quyết triệt để, cung cấp cái nhìn sâu sắc từ nền tảng toán học của quy hoạch động Viterbi đến vi kiến trúc phần cứng và kỹ thuật xử lý tín hiệu thực nghiệm.}
